% Options for packages loaded elsewhere
\PassOptionsToPackage{unicode}{hyperref}
\PassOptionsToPackage{hyphens}{url}
\PassOptionsToPackage{dvipsnames,svgnames,x11names}{xcolor}
%
\documentclass[
  10pt,
  a4paper,
]{article}
\usepackage{amsmath,amssymb}
\usepackage{iftex}
\ifPDFTeX
  \usepackage[T1]{fontenc}
  \usepackage[utf8]{inputenc}
  \usepackage{textcomp} % provide euro and other symbols
\else % if luatex or xetex
  \usepackage{unicode-math} % this also loads fontspec
  \defaultfontfeatures{Scale=MatchLowercase}
  \defaultfontfeatures[\rmfamily]{Ligatures=TeX,Scale=1}
\fi
\usepackage{lmodern}
\ifPDFTeX\else
  % xetex/luatex font selection
\fi
% Use upquote if available, for straight quotes in verbatim environments
\IfFileExists{upquote.sty}{\usepackage{upquote}}{}
\IfFileExists{microtype.sty}{% use microtype if available
  \usepackage[]{microtype}
  \UseMicrotypeSet[protrusion]{basicmath} % disable protrusion for tt fonts
}{}
\makeatletter
\@ifundefined{KOMAClassName}{% if non-KOMA class
  \IfFileExists{parskip.sty}{%
    \usepackage{parskip}
  }{% else
    \setlength{\parindent}{0pt}
    \setlength{\parskip}{6pt plus 2pt minus 1pt}}
}{% if KOMA class
  \KOMAoptions{parskip=half}}
\makeatother
\usepackage{xcolor}
\usepackage[margin=21mm]{geometry}
\usepackage{longtable,booktabs,array}
\usepackage{calc} % for calculating minipage widths
% Correct order of tables after \paragraph or \subparagraph
\usepackage{etoolbox}
\makeatletter
\patchcmd\longtable{\par}{\if@noskipsec\mbox{}\fi\par}{}{}
\makeatother
% Allow footnotes in longtable head/foot
\IfFileExists{footnotehyper.sty}{\usepackage{footnotehyper}}{\usepackage{footnote}}
\makesavenoteenv{longtable}
\setlength{\emergencystretch}{3em} % prevent overfull lines
\providecommand{\tightlist}{%
  \setlength{\itemsep}{0pt}\setlength{\parskip}{0pt}}
\setcounter{secnumdepth}{-\maxdimen} % remove section numbering
\usepackage{fontspec}
\setmainfont{NotoSansCJKsc-Regular.otf}[Path=fonts/,BoldFont=NotoSansCJKsc-Bold.otf,ItalicFont=NotoSansCJKsc-Regular.otf,BoldItalicFont=NotoSansCJKsc-Bold.otf]
\setsansfont{NotoSansCJKsc-Regular.otf}[Path=fonts/,BoldFont=NotoSansCJKsc-Bold.otf]
\setmonofont{DejaVu Sans Mono}[Scale=0.82]
\XeTeXlinebreaklocale "zh"
\XeTeXlinebreakskip=0pt plus 1pt
\usepackage{amsmath,amssymb}
\usepackage{fancyhdr}
\usepackage{etoolbox}
\pagestyle{fancy}
\fancyhf{}
\fancyhead[L]{\footnotesize Research W-GQ-R2-T}
\fancyhead[R]{\footnotesize 目标充分 Maxwell 推断}
\fancyfoot[C]{\footnotesize \thepage}
\setlength{\headheight}{16pt}
\setlength{\parindent}{0pt}
\setlength{\parskip}{5pt}
\emergencystretch=3em
\allowdisplaybreaks
\renewcommand{\arraystretch}{1.18}
\AtBeginEnvironment{longtable}{\footnotesize}
\AtBeginDocument{\hypersetup{colorlinks=true,linkcolor=blue!45!black,urlcolor=blue!55!black}}
\providecommand{\tightlist}{\setlength{\itemsep}{1pt}\setlength{\parskip}{1pt}}


\AtBeginDocument{\renewcommand{\contentsname}{目录}}
\let\originaltableofcontents\tableofcontents
\renewcommand{\tableofcontents}{\originaltableofcontents\clearpage}
\AtBeginDocument{\Urlmuskip=0mu plus 1mu}
\ifLuaTeX
  \usepackage{selnolig}  % disable illegal ligatures
\fi
\IfFileExists{bookmark.sty}{\usepackage{bookmark}}{\usepackage{hyperref}}
\IfFileExists{xurl.sty}{\usepackage{xurl}}{} % add URL line breaks if available
\urlstyle{same}
\hypersetup{
  colorlinks=true,
  linkcolor={Maroon},
  filecolor={Maroon},
  citecolor={Blue},
  urlcolor={Blue},
  pdfcreator={LaTeX via pandoc}}

\author{}
\date{}

\begin{document}

\begin{titlepage}
\thispagestyle{empty}
\vspace*{18mm}
{\color{blue!45!black}\Large Research W-GQ-R2-T\par}
\vspace{12mm}
{\LARGE\bfseries 当墙体材料不可辨识时，\\[5mm]目标究竟需要哪些传播信息？\par}
\vspace{12mm}
{\large Target-Sufficient Maxwell Inference\\[3mm]Beyond Wall Reconstruction\par}
\vspace{16mm}
{\large 中文深度理论研究报告\par}
\vspace{5mm}
{2026年10月29日\par}
\vfill
\noindent\rule{\textwidth}{0.6pt}
\vspace{4mm}
\textbf{核心问题}\par
在普通二维非磁、非均匀墙体与完整多次散射条件下，能否不准确恢复墙体材料，而可靠地恢复墙后目标？\par
\vspace{5mm}
\textbf{交付内容}\par
固定证据核对；条件定理、证明与反例；目标充分传播和概率计算；近邻文献；E1/E2/E3实验合同；Codex执行协议。\par
\vspace{5mm}
\textbf{证据边界}\par
本轮为理论研究与小规模代数核对，没有新增PDE实测。候选算法的目标准确性和实际净加速仍须由冻结合同下的全波实验确认。\par
\vspace{10mm}
\end{titlepage}

{
\hypersetup{linkcolor=}
\setcounter{tocdepth}{2}
\tableofcontents
}
\textbf{研究范围：} 普通二维、被动、非磁、\(E_z\)
scalar-TM；静态未知墙与未知目标；完整非线性材料散射；有限源、有限接收器、有限频带；目标可辨识性、目标概率可靠性及实际物理计算。

\textbf{证据状态：}
已读取四个指定固定提交的报告、BGGA最终机器判决与成本账本，以及上一轮35页W-GQ报告。本轮没有新增PDE运行。文中``命题''表示在所列假设下给出的推导；``候选贡献''表示尚需文献定位与决定性全波实验支持的论文主张。小型矩阵微分核对只检验代数，不作为Maxwell实验证据。

\hypertarget{ux7ed3ux8bbaux4e0eux672cux8f6eux771fux6b63ux7684ux63a8ux8fdb}{%
\section{结论与本轮真正的推进}\label{ux7ed3ux8bbaux4e0eux672cux8f6eux771fux6b63ux7684ux63a8ux8fdb}}

\textbf{可以不准确恢复墙体材料而准确恢复某些隐藏目标任务；普通二维非均匀非磁介质也允许这种现象。但它既不是任意墙类下的普遍保证，也不是把墙改写为Green后自动获得的性质。}

本轮建议把问题具体化为：在给定噪声、频带、几何间隔、目标族和材料幅度约束下，哪些墙传播变化足以把不同目标相互补偿？必须辨认或概率保留的是这些变化；其余材料自由度可以留在等效类内。

相较上一轮，本报告推进四项内容。

\begin{enumerate}
\def\labelenumi{\arabic{enumi}.}
\tightlist
\item
  \textbf{二维物理正例与反例同时成立。}
  有界细微结构具有相同局部均值时，可以保持\(O(1)\)材料差而使整个声明目标族的完整响应差为\(O(\eta)\)。另一方面，若合法墙扰动的全波产品核满射，隐函数定理给出真实的非线性墙--目标精确补偿。两个结果分别回答有限分辨率和无限小尺度问题。
\item
  \textbf{压缩对象改为墙返回传播。}
  写\(C_W=C_0+H_W\)，保留含奇异核的目标自散射，只压缩经过墙--目标间隔的返回\(H_W\)。可推导频率、间距和闭环放大显式出现的误差界。这是标准散射分解的新任务用途，不把分解本身作为原创。
\item
  \textbf{似然对比需要``目标增量＋背景残差投影''。}
  目标增量相近仍不足够。完整推导给出强反馈下四块传播的对比余切，特别是成对前向电流--残差伴随电流外积，而非只看Green能量或单点Jacobian弱模。
\item
  \textbf{背景质量只在目标结论不同的时候需要精算。}
  如果各背景条件目标posterior已经一致，准确的背景posterior权重完全无关；如果仅共同决策一致，只能认证该决策。其一般概率基础是混合分布的收缩性质；新算法成败取决于能否用便宜的Maxwell证书判定这种一致，而不遍历所有原子。
\end{enumerate}

\textbf{优先论文主线：A，目标与墙材料的分辨率分离。}
B作为物理表示，C作为须通过冷启动成本检验的算法增量。单纯``任务充分Green''``无需完整墙图''``风险驱动ROM''都已有强近邻。较有希望的单一主张是：

\begin{quote}
\textbf{有限孔径下目标稳定的Maxwell背景传播类：用完整墙--目标返回作用界定目标分辨率，并只为会改变目标结论的传播不确定性付费。}
\end{quote}

目前对一般未知墙系统、一般目标posterior和实际净加速均不作成功宣告。

\hypertarget{ux56faux5b9aux8bc1ux636eux91cdux65b0ux6574ux7406}{%
\section{1.
固定证据重新整理}\label{ux56faux5b9aux8bc1ux636eux91cdux65b0ux6574ux7406}}

\hypertarget{ux5df2ux6838ux5bf9ux7684ux7248ux672c}{%
\subsection{1.1
已核对的版本}\label{ux5df2ux6838ux5bf9ux7684ux7248ux672c}}

\begin{itemize}
\tightlist
\item
  {[}E1{]} BGGA G0/G1：73ba6eec固定报告、GATE\_DECISION和COST\_SUMMARY。
\item
  {[}E2{]}
  BGGA第一版：6fd50478。它记录实现与资格前检查，新增全波实验为零；不能作为G0的第二份独立物理重复。
\item
  {[}E3{]} 旧Research W：958a4787的P0-R至P1报告。
\item
  {[}E4{]} W-SLAM：52351aee报告。
\item
  {[}E5{]}
  上一轮Research\_W\_GQ\_Report\_ZH.pdf，35页，报告日期2026年10月28日。
\end{itemize}

本轮没有拿到新的Codex
W-GQ实验账本。因此，旧GQ的数学方案仍按理论与提案处理，没有补写新的实测结果。完整链接见文末。

\hypertarget{ux54eaux4e9bux8bc1ux636eux5e94ux6210ux4e3aux57faux7840}{%
\subsection{1.2
哪些证据应成为基础}\label{ux54eaux4e9bux8bc1ux636eux5e94ux6210ux4e3aux57faux7840}}

\begin{longtable}[]{@{}
  >{\raggedright\arraybackslash}p{(\columnwidth - 4\tabcolsep) * \real{0.3333}}
  >{\raggedright\arraybackslash}p{(\columnwidth - 4\tabcolsep) * \real{0.3333}}
  >{\raggedright\arraybackslash}p{(\columnwidth - 4\tabcolsep) * \real{0.3333}}@{}}
\toprule\noalign{}
\begin{minipage}[b]{\linewidth}\raggedright
证据
\end{minipage} & \begin{minipage}[b]{\linewidth}\raggedright
真实支持
\end{minipage} & \begin{minipage}[b]{\linewidth}\raggedright
本轮如何使用
\end{minipage} \\
\midrule\noalign{}
\endhead
\bottomrule\noalign{}
\endlastfoot
G0，28条精确算子对照 &
同公共FEM空间最大内部场相对误差约\(2.04\times10^{-13}\)；材料依赖右端与全反馈保留
& 停止重复证明标准消元是否可实现 \\
12组独立离散比较 & 白化差0.006925--0.020328 &
作为已测配置资格，不作连续或全prior统一证书 \\
C2/C3 & 墙材料不确定性与窄目标任务变化可以分离 &
追查分离机制，不宣布一般目标恢复 \\
旧W联合补偿 &
\(d_{\rm joint}=0.798716\)，目标移位\(0.1\lambda_0\)，\(d_{\rm cal}=3.560512\)
& 保留为独立的条件全波补偿见证 \\
G1 M1/M3 & joint TV均值0.113268，联合门13/24，位置regret门24/24 &
证明决策保真可以弱于joint保真 \\
W-SLAM有限joint & 普通incremental已等价于有限batch reference &
普通联合推断是强基线，不以永久冻结作唯一对手 \\
\end{longtable}

C2的相对信息保留约0.60--0.95，但有限posterior仍保留98.0\%--98.9\%墙先验方差。C3保留约99.73\%--99.92\%，窄固定目标任务差约0.040--0.052。C2与C3复用相同三个parent，12个fixture只有9个独立物理世界；局部nuisance合同与有限prior也不同。{[}E1{]}

C4\_1保持目标位置不变，只改变墙介电参数。单复端口距离0.729007，主两发16收约2.7699，25个冻结目标probe最大差3.068166。它说明任务扩展会破坏窄观测等价，不能写成新增完整孔径目标位移多解。{[}E1{]}

\hypertarget{ux54eaux4e9bux8defux7ebfux5e94ux505cux6b62ux54eaux4e9bux95eeux9898ux4ecdux5f00ux653e}{%
\subsection{1.3
哪些路线应停止，哪些问题仍开放}\label{ux54eaux4e9bux8defux7ebfux5e94ux505cux6b62ux54eaux4e9bux95eeux9898ux4ecdux5f00ux653e}}

\begin{longtable}[]{@{}
  >{\raggedright\arraybackslash}p{(\columnwidth - 2\tabcolsep) * \real{0.5000}}
  >{\raggedright\arraybackslash}p{(\columnwidth - 2\tabcolsep) * \real{0.5000}}@{}}
\toprule\noalign{}
\begin{minipage}[b]{\linewidth}\raggedright
当前实现
\end{minipage} & \begin{minipage}[b]{\linewidth}\raggedright
裁决及理由
\end{minipage} \\
\midrule\noalign{}
\endhead
\bottomrule\noalign{}
\endlastfoot
118墙组全部构建2257端口Green & 停止作为默认部署方案；setup
RHS和存储过大 \\
默认\([0,1]\)未知likelihood包络M4 &
不再扩大相同方案；全部24例仍查完2950原子 \\
要求完整current/future响应接近的M5 &
不再无机制地扩大；118/118组保留，没有merge \\
angle-only选墙M2 & 拒绝该普遍规则；不否定局部主角度理论 \\
固定22列goal-ROM & 现有held-out精度失败；不代表所有返回传播降阶都失败 \\
``小regret＝目标准确'' &
拒绝；完整参考自身位置RMSE约\(0.595\lambda_0\) \\
``预计算bank后查表很快＝PDE加速'' & 拒绝；G1实际安全避免的物理调用为0 \\
\end{longtable}

同进程query kernel：Full 0.33988秒，Woodbury 0.57170秒，Schur
0.87109秒，背景LU预条件Krylov
0.47185秒。22列goal-ROM约0.01114秒，但五个held-out白化误差0.237、0.837、5.331、4.249、6.517，均未达0.1。端口setup另付2257
RHS、7.64 CPU秒；Green images约1.026 GB。{[}E1{]}

\textbf{账本优先于报告内快照。} 最终COST\_SUMMARY记录新物理3934 RHS、482
FE分解，wrapper
CPU合计435.690883秒；正文早期421.3秒不是最终总额。旧bank已付11800
RHS、2950 FE、9288.02304 CPU秒。没有完整冷部署加速证明。{[}E1{]}

M0真实位置MSE约\(0.35454\lambda_0^2\)，M1/M3约\(0.35496\lambda_0^2\)。所以新方法即使完美复制M0，也仍未完成``准确目标''的科学目标。

\hypertarget{ux6570ux5b66ux5408ux540cux5148ux628aux4e09ux79cdux95eeux9898ux5206ux5f00}{%
\section{2.
数学合同：先把三种问题分开}\label{ux6570ux5b66ux5408ux540cux5148ux628aux4e09ux79cdux95eeux9898ux5206ux5f00}}

统一实化、白化观测为 \[
y=F(W,X,\gamma)+n,\qquad n\sim N(0,I_m).
\]
\(W,X,\gamma\)在所有源、频率间共享；不同目标世界不是运动轨迹。主合同两发16收为32个复观测、64个实坐标。现有每个实虚坐标\(\sigma=0.01\)，不除通道数平方根。

两无噪声均值的Mahalanobis距离\(d\)，与真值对含噪观测的残差范数不是同一阈值。后者的标准噪声半径为
\[
\delta_{m,\alpha}=\sqrt{\chi^2_{m,1-\alpha}}.
\]
新增具有非零条件信息的实际测量才可能增加信息；换Green、算伴随、读取缓存不增加信息。

必须分别报告：

\begin{itemize}
\tightlist
\item
  \textbf{数据可辨识性：} 观测是否排除不同目标；
\item
  \textbf{先验条件下推断：} 哪些解释在proper prior下有多少质量；
\item
  \textbf{计算逼近：} 算法是否忠实于前两者定义的目标。
\end{itemize}

强prior可以改善Bayes风险，却不证明数据本身可辨。profile最优残差可以证明某个解释存在，不能给其posterior质量。

目标任务\(\tau(X)\)须具体。位置、尺寸、介电率、损耗和形状分支具有不同信息需求。连续参数的Jacobian不能代表跨形状分支的可辨性；后者需有限secant或概率事件。

\hypertarget{ux6700ux5c11ux4f20ux64adux4fe1ux606fux6709ux4e09ux4e2aux4e0dux540cux7b54ux6848}{%
\section{3.
``最少传播信息''有三个不同答案}\label{ux6700ux5c11ux4f20ux64adux4fe1ux606fux6709ux4e09ux4e2aux4e0dux540cux7b54ux6848}}

\hypertarget{ux5bf9ux6574ux4e2aux672aux77e5ux76eeux6807ux65cfux7684ux7269ux7406ux5145ux5206ux6027}{%
\subsection{3.1
对整个未知目标族的物理充分性}\label{ux5bf9ux6574ux4e2aux672aux77e5ux76eeux6807ux65cfux7684ux7269ux7406ux5145ux5206ux6027}}

物理充分表示必须保留 \[
X\longmapsto F(W,X)
\]
在全部声明目标族上的作用，允许内部状态坐标变换，但不允许改变物理目标标签。它不必保存墙材料图，也不必保存完整2257列Green。

一套足够的表示是：目标可达且可观测的边界照明、接收、返回传播，以及目标自身散射。若目标族非常丰富，这套表示可能仍然很大。``目标只有二维位置''不推出``两列Green足够''。

\hypertarget{ux5bf9ux5f53ux524dux6570ux636eux7684ux76eeux6807posteriorux5145ux5206ux6027}{%
\subsection{3.2
对当前数据的目标posterior充分性}\label{ux5bf9ux5f53ux524dux6570ux636eux7684ux76eeux6807posteriorux5145ux5206ux6027}}

在固定prior与噪声模型下，目标posterior只需要边缘likelihood的正比例类 \[
[m_y(X)]=\{c\,m_y(X):c>0\}.
\]
若任务仅为\(t=\tau(X)\)，还可先对同一\(t\)内的其他目标变量积分。这个数学最小对象比全传播更弱，但它本身不是便宜算法：计算\(m_y\)仍可能很贵。

当前\(y\)上的统计等价不自动适用于另一噪声实现、更大孔径或新目标族。

\hypertarget{ux5bf9ux56faux5b9aux51b3ux7b56ux7684ux5145ux5206ux6027}{%
\subsection{3.3
对固定决策的充分性}\label{ux5bf9ux56faux5b9aux51b3ux7b56ux7684ux5145ux5206ux6027}}

平方位置损失下，Bayes动作只需目标位置均值；完整可信域需要更多信息。两个posterior可均值相同而方差、模式、覆盖率完全不同。

因此没有一个与任务无关的``最小Green维数''。本研究应报告 \[
\mathcal C_{\rm prop}(\epsilon;\mathcal T,\mathcal A),\quad
\mathcal C_{\rm post}(\epsilon;y,\pi_0),\quad
\mathcal C_{\rm act}(\epsilon;y,\ell)
\] 三个任务相关复杂度，而不是把所有问题压成一个端口秩。

\hypertarget{ux547dux9898a1ux666eux901aux4e8cux7ef4ux5faeux7ed3ux6784ux5899ux53efux4ee5ux6750ux6599ux5deeux5927ux5168ux76eeux6807ux65cfux54cdux5e94ux5deeux5c0f}{%
\section{4.
命题A1：普通二维微结构墙可以材料差大、全目标族响应差小}\label{ux547dux9898a1ux666eux901aux4e8cux7ef4ux5faeux7ed3ux6784ux5899ux53efux4ee5ux6750ux6599ux5deeux5927ux5168ux76eeux6807ux65cfux54cdux5e94ux5deeux5c0f}}

\hypertarget{ux7269ux7406ux6761ux4ef6}{%
\subsection{4.1 物理条件}\label{ux7269ux7406ux6761ux4ef6}}

采用非磁\(E_z\)模型 \[
-\Delta u-k^2\varepsilon_r u=f.
\]
介电率位于零阶项。这一点不能直接推广至介电率进入主部的另一极化，也不能直接推广至三维矢量Maxwell。
定理在包含墙和目标支撑的有界Lipschitz域\(\Omega\)上表述，外边界采用精确辐射DtN条件；令\(\mathcal H=H^1(\Omega)\)，载荷空间为其对偶\(\mathcal H\prime\)。实际PML/FEM误差另计，不能以公共离散恒等式替代该连续合同。

设 \[
\varepsilon_{W,\eta}^{(j)}(r)=\bar\varepsilon_W(r)+a(r)p_j(r/\eta),\quad j=1,2,
\]
其中\(p_j\in L^\infty_{\rm per}\)为二维周期零均值函数，\(a\)先取墙内光滑紧支撑，并在正测度区域非零；材料实部为正、损耗非负、对比度不随\(\eta\)发散。
频率相关时只能使用跨频共享的已声明被动色散参数，微结构与原函数界在频带上统一；不能逐频独立拟合。不同微结构可以是颗粒、二维交错条带、孔隙；它们共享局部均值，但不是同一材料图。

若\(p_1-p_2\)非零，则 \[
\|\varepsilon_{W,\eta}^{(1)}-\varepsilon_{W,\eta}^{(2)}\|_{L^2(D_W)}
\longrightarrow c_a>0.
\]

目标\(X\)可遍历一个有界族，位置、尺寸、材料、损耗及形状分支均可变。要求平均背景下完整联合算子满足
\[
\sup_{X\in\mathcal T,k\in\mathcal K}
\|\bar{\mathcal A}_{X,k}^{-1}\|_{\mathcal H\prime\to\mathcal H}\le M.
\]
这是全波稳定裕度，允许强目标散射，不要求Born级数收敛。源、接收采用有界端口泛函；理想线源应减去已知奇异入射场，远离材料的点接收再使用局部正则性。

\hypertarget{ux7ed3ux8bba}{%
\subsection{4.2 结论}\label{ux7ed3ux8bba}}

对足够小\(\eta\)， \[
\boxed{
\sup_{X,k}\|F(W_\eta^{(1)},X;k)-F(W_\eta^{(2)},X;k)\|
\le C k_{\max}^2M^2\eta .
}
\]
常数包含源、白化观测与微结构包络范数。该界同时控制声明族中全部未知目标，保留墙--目标所有阶反馈。

这证明的是\textbf{有限频带、有限噪声尺度上的近似传播等价}，不是每个有限\(\eta\)下的精确等值材料族。

\hypertarget{ux8bc1ux660e}{%
\subsection{4.3 证明}\label{ux8bc1ux660e}}

零均值允许构造有界周期向量场\(\Psi_j\)，使\(\nabla_y\cdot\Psi_j=p_j\)。于是
\[
\int a\,p_j(r/\eta)u\bar v\,dr
=-\eta\int\Psi_j(r/\eta)\cdot\nabla(a u\bar v)\,dr .
\] 由Cauchy--Schwarz， \[
\|M_{a p_j(\cdot/\eta)}\|_{\mathcal H\to\mathcal H\prime}\le C\eta.
\]
写\(\mathcal A_\eta=\bar{\mathcal A}-k^2M_{ap_j}\)。若\(Ck_{\max}^2M\eta<1/2\)，扰动引理给出\(\|\mathcal A_\eta^{-1}\|\le2M\)，且
\[
u_\eta-\bar u
=\mathcal A_\eta^{-1}k^2M_{ap_j}\bar{\mathcal A}^{-1}f.
\]
所以\(\|u_\eta-\bar u\|_{H^1}\le2Ck_{\max}^2M^2\eta\|f\|\)。对两种微结构与平均模型比较并作用观测泛函，即得结论。

没有截断联合逆算子，因此没有删除强散射。若微结构切到墙边界，需加入边界迹与修正项；一般总体仍有\(O(\eta)\)项，不能随意宣称\(O(\eta^2)\)。{[}R2{]}

\hypertarget{ux4eceux4f20ux64adux7b49ux4ef7ux5230ux6b63ux786eux76eeux6807}{%
\subsection{4.4
从传播等价到正确目标}\label{ux4eceux4f20ux64adux7b49ux4ef7ux5230ux6b63ux786eux76eeux6807}}

若平均模型在声明区域满足 \[
\|\bar F(\zeta_1,X_1)-\bar F(\zeta_2,X_2)\|
\ge c_X d_{\mathcal T}(X_1,X_2),
\] 且每个微结构世界对平均模型的响应误差\(\le\epsilon_\eta\)，则 \[
\boxed{
\omega_{\mathcal T}(\delta)
\le(\delta+2\epsilon_\eta)/c_X.
}
\]

对于实际算法，若输出\((\widehat\zeta,\widehat X)\)满足有效模型拟合残差\(\le\epsilon_{\rm fit}\)，则在\(\|n\|\le\delta_n\)时
\[
d_{\mathcal T}(\widehat X,X^\star)
\le(\epsilon_{\rm fit}+\delta_n+\epsilon_\eta)/c_X.
\] 有效墙与真实微结构墙的材料差仍可为\(O(1)\)。

\textbf{这里没有免费假设\(c_X\)。}
下一节给出可检验的局部到有限尺度条件；跨分支还需独立secant证据。该命题不覆盖任意宏观错误墙、任意大空腔或所有目标位置。

\hypertarget{ux7269ux7406ux89e3ux91caux4e0eux9002ux7528ux8fb9ux754c}{%
\subsection{4.5
物理解释与适用边界}\label{ux7269ux7406ux89e3ux91caux4e0eux9002ux7528ux8fb9ux754c}}

在该极化下，有界小尺度介电变化作为零阶势进入方程；其对有限频带波场的作用经过积分平滑。大材料误差与小传播误差因此可以共存。提高频率会增加界中的频率因子；提高对比度到微共振尺度可能超出有界对比度假设；接近联合共振会放大\(M\)。目标靠近墙也可能使整个目标族的统一\(M\)恶化，但靠近本身不自动违反A1。是否失效，应检查同一合同下\(Ck_{\max}^2M\eta\)与最终白化误差是否仍足够小。

均匀化本身已有成熟理论。本轮候选增量是``对未知目标族统一＋完整反馈＋有限目标分辨率''的组合与其可验证常数，而非重新命名有效介质。{[}R2{]}

\hypertarget{ux547dux9898a2ux4e0eux53cdux547dux9898ux76eeux6807ux4f55ux65f6ux80fdux6bd4ux5899ux4f53ux66f4ux53efux8fa8}{%
\section{5.
命题A2与反命题：目标何时能比墙体更可辨？}\label{ux547dux9898a2ux4e0eux53cdux547dux9898ux76eeux6807ux4f55ux65f6ux80fdux6bd4ux5899ux4f53ux66f4ux53efux8fa8}}

\hypertarget{ux5b8cux6574ux5168ux6ce2ux7684ux5bf9ux79f0ux6027ux6b63ux4f8b}{%
\subsection{5.1
完整全波的对称性正例}\label{ux5b8cux6574ux5168ux6ce2ux7684ux5bf9ux79f0ux6027ux6b63ux4f8b}}

允许墙体在二维内任意非均匀，但合法墙类关于\(x\mapsto-x\)镜像对称；可以有对称空腔与不规则边界。源接收成镜像对，白化保留该对称性。令白化后的数据镜像算子满足\(\mathsf P^2=I\)与\(\mathsf P^*=\mathsf P\)，因此\(P_\pm=(I\pm\mathsf P)/2\)是正交投影。

对镜像对称目标的横向平移\(x\)，PDE等变性给出 \[
F(W,-x)=\mathsf P F(W,x).
\] 在\(x=0\)，合法墙扰动落偶子空间，横向位移导数落奇子空间： \[
P_-A=0,\qquad P_+B_x=0.
\] 若仪器nuisance也只占偶子空间，且\(B_x\ne0\)， \[
\boxed{
K_{x\mid W,\gamma}=\|B_x\|^2>0.
}
\]

证明仅用算子对称性与正交分解，不使用Born近似。墙有无限多材料自由度，有限数据仍不能恢复全部墙；横向目标在此点却保留全部局部信息。

若对整个墙类\(\|b_W\|\ge\beta>0\)且三阶导数界为\(L_3\)， \[
P_-F(W,x)=x b_W+O(x^3),
\] 故 \[
\|P_-F(W,x)\|\ge
|x|\left(\beta-\frac{L_3r^2}{6}\right),\quad |x|\le r.
\]
当选择\(r\)使\(\beta-L_3r^2/6>0\)时，它给出相对居中真值的有限定位界。若该系数非正，上式只是无效的下界，不能据此认证定位。

\textbf{边界：}
不自动给任意\(x_1,x_2\)间的统一稳定性，因为\(b_W\)随墙变化；尺寸、损耗等偶任务也不受该奇偶保护。非对称墙块、未知左右增益或横向天线偏移都可破坏结论。该正例适合作机制对照，而非整个成像系统的成功证明。

\hypertarget{ux7cbeux786eux975eux7ebfux6027ux53cdux4f8bux5408ux6cd5ux5899ux53efux4ee5ux5438ux6536ux76eeux6807ux53d8ux5316}{%
\subsection{5.2
精确非线性反例：合法墙可以吸收目标变化}\label{ux7cbeux786eux975eux7ebfux6027ux53cdux4f8bux5408ux6cd5ux5899ux53efux4ee5ux5438ux6536ux76eeux6807ux53d8ux5316}}

实化后的全波墙导数具有产品核 \[
D_WF[h]_{rs}=k^2\int_{D_W}h(r)u_s(r)v_r(r)\,dr,
\] \(v_r\)为相应互易或伴随场，采用与实化内积一致的约定。

假设在材料约束内部，存在\(m\)个允许的墙内实参数方向\(\phi_j\)，使 \[
\mathsf A_{\ell j}=D_WF[\phi_j]_\ell
\] 可逆。令\(W(c)=W_0+\sum_jc_j\phi_j\)。对 \[
H(c,\theta)=F(W(c),X(\theta))-F(W_0,X(\theta_0))
\] 应用隐函数定理，存在\(c(\theta)\)使 \[
\boxed{F(W(c(\theta)),X(\theta))=F(W_0,X(\theta_0)).}
\]

这是完整非线性全波等值族。墙内介电率与损耗的足够小正负增量仍可保持普通被动非磁条件。\textbf{满射导致的歧义不能靠再看二阶或三阶消除。}

该条件有普通二维构造：近弱对比度时，平面波导数核趋于\(e^{ik(d_s-\widehat r)\cdot r}\)。不同独立散射向量的指数函数在墙内开放集上线性独立，可选有限bump中心使积分矩阵非奇异。加入小而严格正的材料和损耗，真实全波导数连续保持满秩，再对真实非线性映射应用隐函数定理。Born核这里只用于证明某个全波正则点存在，不用于最终前向近似。

多频反例必须在共同色散模型中验证满射；不能逐频自由造墙。有限天线阵列还必须去掉互易冗余观测。

\hypertarget{ux4e3aux4ec0ux4e48ux5b83ux4e0eux5faeux7ed3ux6784ux6b63ux4f8bux4e0dux77dbux76fe}{%
\subsection{5.3
为什么它与微结构正例不矛盾？}\label{ux4e3aux4ec0ux4e48ux5b83ux4e0eux5faeux7ed3ux6784ux6b63ux4f8bux4e0dux77dbux76fe}}

精确投影\(P_A\)只看响应方向，不看该方向需要多大材料变化。一个极小但非零的墙奇异值仍进入\(\operatorname{Ran}A\)。无限小目标位移可以被微小合法墙变化补偿，而有限目标位移可能要求超出材料幅度范围的墙变化。

因此可能同时有：

\begin{itemize}
\tightlist
\item
  精确局部\(K_{X\mid W}=0\)；
\item
  物理幅度限制下\(\omega_X(\delta)\le C(\delta+\eta)\)；
\item
  有限噪声下目标达到某个实用分辨率。
\end{itemize}

截掉小墙奇异值后的Schur必须标为``噪声/幅度约束下的有效信息''，不能冒充纯数据精确可辨识性。这也是本轮比单点谱分析更重要的尺度区分。

\hypertarget{ux6709ux9650ux975eux7ebfux6027ux7a33ux5b9aux6027ux4e0eux9ad8ux9636ux8865ux507f}{%
\section{6.
有限非线性稳定性与高阶补偿}\label{ux6709ux9650ux975eux7ebfux6027ux7a33ux5b9aux6027ux4e0eux9ad8ux9636ux8865ux507f}}

\hypertarget{ux53efux8ba1ux7b97ux7684ux5c40ux90e8ux5230ux6709ux9650ux5c3aux5ea6ux6761ux4ef6}{%
\subsection{6.1
可计算的局部到有限尺度条件}\label{ux53efux8ba1ux7b97ux7684ux5c40ux90e8ux5230ux6709ux9650ux5c3aux5ea6ux6761ux4ef6}}

对有限维宏观参数\(z=(\zeta,\theta,\gamma)\)，先固定无害gauge。设\(J_0=D\bar F(z_0)\)满列秩，并在凸参数盒\(\mathcal B\)内验证
\[
q=\sup_{z\in\mathcal B}
\|J_0^\dagger(D\bar F(z)-J_0)\|<1.
\] 沿两点间线段积分， \[
J_0^\dagger(\bar F(z_1)-\bar F(z_2))
=(z_1-z_2)+e,\quad \|e\|\le q\|z_1-z_2\|.
\] 故 \[
\boxed{
\|z_1-z_2\|
\le\frac{\|J_0^\dagger\|}{1-q}
\|\bar F(z_1)-\bar F(z_2)\|.
}
\]

这是一个实际可用的有限盒充分条件。它要求可辨的是宏观传播参数与目标，不要求微结构材料可辨。它仍可能保守；若宏观参数有无害核，应先采用合法局部商坐标，不能直接给不满列矩阵套伪逆结论。

对仿射材料算子\(\mathcal A=\mathcal A_0-\sum_i z_iM_i\)、\(u=\mathcal A^{-1}f\)：
\[
D_iF=O\mathcal A^{-1}M_i u,
\] \[
D_{ij}F
=O\mathcal A^{-1}
(M_i\mathcal A^{-1}M_j+M_j\mathcal A^{-1}M_i)u.
\] 若完整逆算子范数\(\le M\)，则 \[
\|D^2F[h,k]\|
\le2\|O\|M^3\|M_h\|\|M_k\|\|f\|.
\]
这给\(q\)与盒半径的联系，也明确显示共振如何缩小稳定邻域。非仿射位置/形状参数需增加\(O\mathcal A^{-1}M_{ij}u\)及载荷导数。

支撑分支之间没有共享光滑盒，必须另求 \[
s_{\rm branch}=\inf_{\theta_1\in\Theta_1,\theta_2\in\Theta_2}
\|F_1(\theta_1)-F_2(\theta_2)\|.
\] 只找到较大的局部最小值不能给这个inf下界。

\hypertarget{ux89c2ux6d4bux51e0ux4f55ux4e0dux662fux7a7aux6cdbux6761ux4ef6}{%
\subsection{6.2
观测几何不是空泛条件}\label{ux89c2ux6d4bux51e0ux4f55ux4e0dux662fux7a7aux6cdbux6761ux4ef6}}

有限维宏观材料与目标导数作为空间分布线性独立时，其Fourier变换不能在一个开放散射向量集合上全部相消。对具有非退化角域的收发几何，可选择有限收发对，使弱散射邻域的参数Jacobian满列。目标位置导数需要非零目标对比度，真空中没有可定位目标。

沿保持解析、可解的材料路径，一个已知非零minor的零点至多离散。因此强散射段可以存在正则点；\textbf{它不给强散射全区统一奇异值下界}。对实际两发16收，只能计算并核验，不能把一般存在性当作该阵列已经合格。

\hypertarget{ux5b8cux6574maxwellux9ad8ux9636ux8865ux507f}{%
\subsection{6.3
完整Maxwell高阶补偿}\label{ux5b8cux6574maxwellux9ad8ux9636ux8865ux507f}}

以下在一个光滑目标分支内讨论，并假设\(\tau\)可作为局部坐标的一部分；离散任务或\(\tau\)的临界点需另作有限secant分析。把墙、仪器和不属于任务的目标参数统一记为\(n\)。设任务方向\(v\)，\(A=D_nF\)、\(B=D_\tau F\)。若\(Aa+Bv=0\)，取
\[
n(t)=n_0+ta+\tfrac12t^2c,\qquad \tau(t)=\tau_0+tv.
\] 则 \[
F(t)-F_0
=\frac{t^2}{2}\{Ac+D^2F[(a,v),(a,v)]\}+O(t^3).
\] 二阶无法由自由nuisance消去的部分为 \[
r_2(a,v)=(I-P_A)D^2F[(a,v),(a,v)].
\] 正确任务量必须对全部合法一阶补偿优化： \[
\chi_2(v)=\inf_{a:\ Aa=-Bv,\ a\ {\rm feasible}}\|r_2(a,v)\|.
\]

一个最小范数\(a\)的曲率大，不证明所有补偿都被排除；\(\ker A\)中的分量可以改变二阶结果。box、passivity与几何边界需使用可行切锥和二阶切集，不能继续用无约束正交投影当精确最优值。

找到残差\(O(t^2)\)的实际曲线，给\(\omega_X(\delta)\gtrsim\sqrt\delta\)型歧义下界。要给反方向上界，需对所有可行曲线给统一正曲率并排除远处与异拓扑分支。

Maxwell混合导数为 \[
D_{WX}F[h,k]
=O\mathcal A^{-1}
(M_h\mathcal A^{-1}M_k+M_k\mathcal A^{-1}M_h)u.
\]
即使墙与目标支撑不重叠、\(M_hM_k=0\)，中间的传播算子也使混合项非零。这正是两种墙--目标散射顺序造成的有限补偿曲率。

\hypertarget{ux65b0ux589eux6d4bux91cfux53efux4ee5ux6539ux5584ux76eeux6807ux800cux4e0dux6062ux590dux5899ux6750ux6599}{%
\section{7.
新增测量可以改善目标而不恢复墙材料}\label{ux65b0ux589eux6d4bux91cfux53efux4ee5ux6539ux5584ux76eeux6807ux800cux4e0dux6062ux590dux5899ux6750ux6599}}

对真正嵌套的观测，且新旧观测使用同一\(n\)， \[
v^*K_{\rm eff}^{\rm new}v
=\inf_h\left\{
\|A_{\rm old}h+B_{\rm old}v\|^2+
\|A_{\rm add}h+B_{\rm add}v\|^2
\right\}
\ge v^*K_{\rm eff}^{\rm old}v.
\]
证明是新增非负项后取同一变量的inf。相关噪声可先做条件创新白化；若增加通道同时削弱原测量SNR，则不属于这个嵌套合同。

新增测量打破目标歧义的具体条件是： \[
\{h:A_{\rm old}h=-B_{\rm old}v\}
\cap
\{h:A_{\rm add}h=-B_{\rm add}v\}
=\varnothing.
\]
它只要求不同观测需要的墙补偿不再相容，不要求墙Jacobian变成单射。剩余墙nullspace仍可很大。

这给出三类设计。

\begin{itemize}
\tightlist
\item
  \textbf{角度互补：} 增加使目标位移导数落入不同产品核的源/接收角度。
\item
  \textbf{共同色散约束下的多频：}
  同一\(\varepsilon_\infty,\ell\)必须同时解释所有频率。逐频自由墙参数会消耗增益。
\item
  \textbf{对称成对测量：} 让特定目标方向进入nuisance无法占据的奇偶通道。
\end{itemize}

对上述嵌套合同，局部efficient信息不减；采用同一proper
prior与最优Bayes决策时，新增数据的期望Bayes风险不增。二者都不保证严格改善，也不保证每一次噪声实现的真值误差下降；近似算法还可能因数值误差偏离该最优性结论。报告须给同parent配对统计，同时保留反常个例。目标loss与wall
loss分别评价；不以墙MSE下降作为选测量唯一目标。

\hypertarget{ux547dux9898bux4fddux7559ux76eeux6807ux81eaux6563ux5c04ux53eaux538bux7f29ux5899ux8fd4ux56de}{%
\section{8.
命题B：保留目标自散射，只压缩墙返回}\label{ux547dux9898bux4fddux7559ux76eeux6807ux81eaux6563ux5c04ux53eaux538bux7f29ux5899ux8fd4ux56de}}

\hypertarget{ux6b63ux786eux5206ux89e3}{%
\subsection{8.1 正确分解}\label{ux6b63ux786eux5206ux89e3}}

在目标域内写 \[
C_W=C_0+H_W,\qquad
T_X^0=V_X(I-C_0V_X)^{-1}.
\]
\(C_0\)保留参考介质的奇异自作用，\(H_W\)只含从目标出发、经墙散射再返回的部分。完整散射为
\[
\boxed{
F(W,X)=d_W+B_WT_X^0(I-H_WT_X^0)^{-1}a_W.
}
\]

这是标准多体散射代数，已有T-matrix、DtN、Schur、Marchenko等强近邻。{[}R3--R8{]}
它的研究价值在于选对降阶对象：高频目标边界电流由\(T_X^0\)处理，不要求它们全部落入一个很小的全局电流子空间。

\hypertarget{ux95f4ux9694ux5e73ux6ed1ux7684ux4e25ux683cux8c31ux754c}{%
\subsection{8.2
间隔平滑的严格谱界}\label{ux95f4ux9694ux5e73ux6ed1ux7684ux4e25ux683cux8c31ux754c}}

设墙与目标之间存在已知均匀空隙\(g>0\)。在平面端口角谱表示中，定义端口返回算子；采用使\(P_g\)为收缩、\(P_K\)为正交投影的共同模态范数：
\[
\mathcal H_W=P_g\mathcal R_WP_g,\qquad
\widehat{P_g u}(\xi)
=e^{ig\sqrt{k^2-\xi^2}}\hat u(\xi).
\]
端口返回\(\mathcal H_W\)与上一节的体电流返回\(H_W\)不属于同一状态空间。继续使用体电流坐标时，应写
\[
H_W=E\mathcal H_WJ,
\]
其中\(J\)将目标电流映到端口出射振幅，\(E\)将端口入射振幅映回目标场；二者由已知参考介质及端口位置定义。也可将\(T_X^0,a_W,B_W\)一并转换到一致端口坐标。

取\(\operatorname{Im}\sqrt{k^2-\xi^2}\ge0\)，假设\(\|\mathcal R_W\|\le M_R\)。这不要求墙横向均匀；\(\mathcal R_W\)可耦合不同模式。

对保留\(|\xi|\le K\)的投影\(P_K\)，\(K>k\)， \[
\boxed{
\|\mathcal H_W-P_K\mathcal H_WP_K\|
\le2M_R e^{-g\sqrt{K^2-k^2}}.
}
\]
证明将误差分成\((I-P_K)P_g\mathcal R_WP_g\)和\(P_KP_g\mathcal R_WP_g(I-P_K)\)，分别估计。

对体电流表示，采用\(\widetilde H_W=E(P_K\mathcal H_WP_K)J\)时，需计入耦合范数：
\[
\|H_W-\widetilde H_W\|
\le2\|E\|\|J\|M_R e^{-g\sqrt{K^2-k^2}}.
\] 此界才可直接代入下一节体电流闭环误差公式。

总体误差只有一个\(g\)的指数，因为不均匀墙可把高模转为低模。只有高--高块才有\(e^{-2g\sqrt{K^2-k^2}}\)。被动性不能直接给全部倏逝模式\(\|\mathcal R_W\|\le1\)。

对横向长度\(L\)、频率为\(2\pi n/L\)的周期合同， \[
r(K)=2\left\lfloor\frac{LK}{2\pi}\right\rfloor+1,
\]
其他波导边界条件应使用其实际横向本征值计数。对上述周期合同，达到\(\epsilon_H\)约需
\[
r=O\left(Lk+\frac Lg\log\frac{M_R}{\epsilon_H}\right).
\] \textbf{开放无限平面的带限空间仍无限维。}
真实有限窗口必须另控截窗、侧边绕射、输入输出迹范数和窗口泄漏；上述计数不能无条件移植为普通开放墙的精确秩。

\hypertarget{ux5f3aux53cdux9988ux8befux5deeux754c}{%
\subsection{8.3
强反馈误差界}\label{ux5f3aux53cdux9988ux8befux5deeux754c}}

令\(T(H)=T_X^0(I-HT_X^0)^{-1}\)。精确恒等式为 \[
T(H)-T(\widetilde H)
=T(H)(H-\widetilde H)T(\widetilde H).
\] 因此 \[
\boxed{
\|F_H-F_{\widetilde H}\|
\le
\|B_WT(H)\|\,
\|H-\widetilde H\|\,
\|T(\widetilde H)a_W\|.
}
\]

上式只改变返回算子，保持\(d,a,B\)不变。比较两个真实背景时，必须同时计入全部四块。令\(\delta d=d_1-d_2\)、\(\delta a=a_1-a_2\)、\(\delta B=B_1-B_2\)、\(\delta H=H_1-H_2\)，则
\[
F_1-F_2=\delta d+\delta B T_1a_1+B_2T_1\delta a+B_2T_1\delta H T_2a_2,
\] 从而 \[
\begin{aligned}
\|F_1-F_2\|\le{}&\|\delta d\|+\|\delta B\|\|T_1a_1\|\\
&+\|B_2T_1\|\bigl(\|\delta a\|+\|\delta H\|\|T_2a_2\|\bigr).
\end{aligned}
\]
仅返回\(H\)相近不能保证不同墙的完整响应相近；完整响应较大也不排除第10节的目标对比相近。

目标内部散射可以很强。危险的是墙--目标闭环放大：若\((I-HT_X^0)^{-1}\)很大，原本很小的返回高模也可能不可忽略。缩小间距、提高频率、改变目标形状或介电率都可能引入新的返回通道，破坏旧等价类。

若\(\widetilde\kappa=\|(I-\widetilde H T_X^0)^{-1}\|\)且
\(\widetilde\kappa\|H-\widetilde H\|\|T_X^0\|<1\)，扰动引理可给完整闭环逆的上界。通常的一范数condition
estimate不能直接充当逆算子范数的经认证上界，也不能自动提供所需的最小奇异值下界。

\hypertarget{ux6700ux5c0fux5b9eux73b0ux5e94ux600eux6837ux8868ux8ff0}{%
\subsection{8.4
最小实现应怎样表述}\label{ux6700ux5c0fux5b9eux73b0ux5e94ux600eux6837ux8868ux8ff0}}

由\(V_X\)与\(C_W\)生成的可达/可观测闭包是安全的充分空间；对有限静态目标参数，它未必是必要的最小空间。可交换参数只给非交换words的对称化Taylor系数，不能直接套free非交换实现唯一性。{[}R15{]}

若目标族能独立激活多数ROI单元，\(\operatorname{span}\{V_Xa\}\)本身就可能满维。二维平移流形也可具有很高的线性包络维数。相比压缩所有内部电流，压缩分离后的墙返回更有物理依据。

一个有利但必须公平使用的特殊结构是：旧2950原子prior只有固定圆盘的25个位置。齐次无限参考介质下，平移不改变孤立目标散射，可用每频一套圆盘Mie/T模板和正确平移算子复用。它来自声明prior，不是使用真值目标；普通T-matrix基线也必须享有同样权利。有限FEM/PML中的混接误差需独立核验。

\hypertarget{ux4e0dux5148ux91cdux5efaux5899ux56feux80fdux5426ux76f4ux63a5ux8bc6ux522bux4f20ux64adux72b6ux6001}{%
\section{9.
不先重建墙图，能否直接识别传播状态？}\label{ux4e0dux5148ux91cdux5efaux5899ux56feux80fdux5426ux76f4ux63a5ux8bc6ux522bux4f20ux64adux72b6ux6001}}

\textbf{有条件能；从一个任意未知场景同时估计任意\(Q\)与\(X\)，一般不能。}
第5节的全波满射反例已经给出信息层面的障碍。

\hypertarget{ux54eaux4e9bgaugeux65e0ux5bb3ux54eaux4e9bux6709ux5bb3}{%
\subsection{9.1
哪些gauge无害，哪些有害？}\label{ux54eaux4e9bgaugeux65e0ux5bb3ux54eaux4e9bux6709ux5bb3}}

\begin{longtable}[]{@{}
  >{\raggedright\arraybackslash}p{(\columnwidth - 2\tabcolsep) * \real{0.5000}}
  >{\raggedright\arraybackslash}p{(\columnwidth - 2\tabcolsep) * \real{0.5000}}@{}}
\toprule\noalign{}
\begin{minipage}[b]{\linewidth}\raggedright
变化
\end{minipage} & \begin{minipage}[b]{\linewidth}\raggedright
对目标任务的影响
\end{minipage} \\
\midrule\noalign{}
\endhead
\bottomrule\noalign{}
\endlastfoot
内部状态相似变换，同时变换\(a,B,C,V_X\)与度量 &
只是同一物理模型的坐标变化，可无害 \\
混合物理空间坐标，却仍把\(V_X\)解释为原位置局部材料 &
会改变目标坐标，不能视作无害 \\
照明幅度与目标对比度互相缩放 & 会污染目标介电率/损耗 \\
未知共享时延与传播光程/目标距离补偿 & 可能保持频率相位而改变深度 \\
未知墙与目标支撑允许重叠、标签没有定义 &
同一总材料场可有不同\(W/X\)标签 \\
任意逐频被动矩阵 & 不自动满足共同因果色散，也不自动可由普通墙实现 \\
\end{longtable}

例如在自由四元组中 \[
V'=sV,\quad C'=C/s,\quad
a'=a/\sqrt{s},\quad B'=B/\sqrt{s}
\]
保持代数响应不变。它证明过宽\(Q\)参数化的危险，不意味着每个变换后的四元组都能由本研究的标量被动墙实现。

若要求固定坐标下完整局部材料算子族不变，允许的内部相似变换受\(SV_X=V_XS\)约束；不能自由旋转潜在空间后继续把坐标当作物理像素。

\hypertarget{ux4e0dux4f7fux7528ux5df2ux77e5ux6821ux51c6ux76eeux6807ux7684ux53efux884cux4fe1ux606fux6765ux6e90}{%
\subsection{9.2
不使用已知校准目标的可行信息来源}\label{ux4e0dux4f7fux7528ux5df2ux77e5ux6821ux51c6ux76eeux6807ux7684ux53efux884cux4fe1ux606fux6765ux6e90}}

\begin{itemize}
\tightlist
\item
  已知源接收几何、相位参考与有限共享仪器模型；
\item
  宏观墙材料类、支撑容器、间隔和共同色散；
\item
  足够互补的多源多频，使目标与允许nuisance响应分离；
\item
  在富集反射数据合同下，用direct arrivals/macromodel识别目标边界波场；
\item
  独立背景或已知负载校准可提供额外信息，但不是本轮默认免费输入。
\end{itemize}

最少额外信息取决于缺失gauge，不能统一说``三个校准目标''或``多一个频率就够''。单纯互易、被动、平滑仍不足以排除目标泄漏。

Marchenko及数据驱动EM
ROM说明``不先获得完整材料图而恢复内部传播''已有严肃先例；它们依赖时域带宽、采样、宏观速度或双侧数据等条件。{[}R3--R6,R9{]}
当前两发16收单频不能自动满足这些合同。

\hypertarget{ux547dux9898cmaxwellux4f3cux7136ux5bf9ux6bd4ux7684ux5145ux5206ux91cfux4e0eux4e25ux683cux8befux5deeux5f0f}{%
\section{10.
命题C：Maxwell似然对比的充分量与严格误差式}\label{ux547dux9898cmaxwellux4f3cux7136ux5bf9ux6bd4ux7684ux5145ux5206ux91cfux4e0eux4e25ux683cux8befux5deeux5f0f}}

\hypertarget{ux76eeux6807ux589eux91cfux4e0eux6b8bux5deeux6295ux5f71ux7f3aux4e00ux4e0dux53ef}{%
\subsection{10.1
目标增量与残差投影缺一不可}\label{ux76eeux6807ux589eux91cfux4e0eux6b8bux5deeux6295ux5f71ux7f3aux4e00ux4e0dux53ef}}

设合法anchor为\(X_0\)，简记 \[
f_{W0}=F(W,X_0),\quad
\Delta_W(X)=F(W,X)-f_{W0},\quad
r_{W0}=y-f_{W0}.
\] 在实白化Gaussian模型中， \[
\boxed{
D_W(X)=
\langle r_{W0},\Delta_W(X)\rangle
-\frac12\|\Delta_W(X)\|^2.
}
\]

若整个声明族的增量都位于观测子空间\(U\)，则背景内部全部目标相对likelihood只需要
\[
\left(P_U r_{W0},\ X\mapsto\Delta_W(X)\right).
\]
背景残差在\(U^\perp\)中的大变化不影响对比。若不同背景具有相同的\(P_Ur_{W0}\)和整个目标增量函数\(X\mapsto\Delta_W(X)\)，则它们的目标likelihood
contrast相同；仅有全部增量位于\(U\)本身并不足够。这给出``完整响应可不同，条件目标推断仍相同''的明确机制。

但\(U\)是整个声明目标--背景族的增量空间，不是从一个真实目标免费得到的子空间。若仅近似满足\(\|(I-P_U)\Delta_W(X)\|\le\rho\)，省略部分的对比误差至多
\[
\|(I-P_U)r_{W0}\|\rho+\tfrac12\rho^2.
\] 高维残差不能从界中删去。

\hypertarget{ux4e24ux80ccux666fux4e4bux95f4ux7684ux7cbeux786eux5deeux5206ux6052ux7b49ux5f0f}{%
\subsection{10.2
两背景之间的精确差分恒等式}\label{ux4e24ux80ccux666fux4e4bux95f4ux7684ux7cbeux786eux5deeux5206ux6052ux7b49ux5f0f}}

记\(\delta\Delta=\Delta_1-\Delta_2\)、\(\delta f_0=f_{10}-f_{20}\)，横线表示两者平均。直接展开得到
\[
\boxed{
D_1(X)-D_2(X)
=
\langle\bar r_X,\delta\Delta\rangle
-\langle\delta f_0,\bar\Delta\rangle.
}
\] 因此 \[
|D_1-D_2|
\le
\|\bar r_X\|\|\delta\Delta\|
+\|P_U\delta f_0\|\|\bar\Delta\|
\]
在\(\bar\Delta\in U\)时成立。完整背景响应可以相差很大，但只要差异近乎垂直于目标增量，并且目标增量本身稳定，contrast仍可接近。

这给出比``完整Green相似度''更具体的比较指标，但还没有保证该指标便宜可算。

\hypertarget{ux5f53ux524dux6570ux636eux7b49ux4ef7ux4e0eux566aux58f0ux91cdux590dux7b49ux4ef7}{%
\subsection{10.3
当前数据等价与噪声重复等价}\label{ux5f53ux524dux6570ux636eux7b49ux4ef7ux4e0eux566aux58f0ux91cdux590dux7b49ux4ef7}}

对数据空间一个开集中的所有\(y\)，两背景contrast完全相同，当且仅当对所有允许\(X\)：
\[
\boxed{
\Delta_1(X)=\Delta_2(X)=\Delta(X),\qquad
\langle f_{10}-f_{20},\Delta(X)\rangle=0.
}
\]
证明：\(D_1-D_2\)是\(y\)的仿射函数；对开集恒零迫使\(y\)系数\(\delta\Delta=0\)，再比较常数项。

所以三个层次必须分开：

\begin{itemize}
\tightlist
\item
  当前\(y\)上的偶然相消；
\item
  对重复噪声稳定的contrast等价；
\item
  对目标、频率、孔径族都稳定的物理传播等价。
\end{itemize}

若比较族事前冻结，\(y=\mu+n\)、\(\|\delta\Delta\|\le e_\Delta\)，均值数据处对比差\(\le b\)，则对\(N\)个比较，以概率至少\(1-\alpha\)，
\[
|D_1-D_2|
\le b+e_\Delta\sqrt{2\log(2N/\alpha)}.
\]
连续目标族需要覆盖数或统一噪声范数界。若方向是看当前噪声后选择的，不能直接套该冻结族结论。

\hypertarget{ux5f3aux53cdux9988ux4e0bux7684ux5bf9ux6bd4ux5bfcux6570ux8ba1ux7b97ux54eaux4e9bux4f20ux64adux65b9ux5411ux624dux6709ux7528}{%
\section{11.
强反馈下的对比导数：计算哪些传播方向才有用？}\label{ux5f3aux53cdux9988ux4e0bux7684ux5bf9ux6bd4ux5bfcux6570ux8ba1ux7b97ux54eaux4e9bux4f20ux64adux65b9ux5411ux624dux6709ux7528}}

\hypertarget{ux56dbux5757ux5bf9ux6bd4ux4f59ux5207}{%
\subsection{11.1
四块对比余切}\label{ux56dbux5757ux5bf9ux6bd4ux4f59ux5207}}

令 \[
T_X=V_X[I-(C_0+H)V_X]^{-1},\quad
q_X=T_Xa,
\] \[
F_X=d+Bq_X,\quad r_X=y-F_X,\quad
p_X=T_X^*B^*r_X.
\] 对\(Q=(d,a,B,H)\)的扰动，有 \[
\delta D_X=\operatorname{Re}\langle\Xi_X,\delta Q\rangle,
\] 其中使用向量或Hilbert--Schmidt内积， \[
\boxed{
\begin{aligned}
\Xi_d&=-(F_X-F_0),\\
\Xi_a&=p_X-p_0,\\
\Xi_B&=r_Xq_X^*-r_0q_0^*,\\
\Xi_H&=p_Xq_X^*-p_0q_0^*.
\end{aligned}
}
\]

证明用\(D_HT_X[E]=T_XET_X\)，故 \[
\delta F=\delta d+\delta Bq_X+BT_X\delta a+BT_X\delta Hq_X,
\] 再对\(\log L_X-\log L_0\)取差。

上式以单源书写。多源时将\(q,p,r\)按源排列，使用Frobenius内积，\(\Xi_H,\Xi_B\)对源求和。若噪声跨源或跨频率相关，须通过白化算子的伴随将likelihood梯度映回物理观测，再按源分块；不能把各源独立白化。

\textbf{物理含义：}
返回误差是否重要，取决于它把哪些实际目标电流送到哪些残差伴随电流，而不是只看\(H\)的能量。对每个频率\(S\)个源，\(\Xi_H\)为两组\(S\)个外积之差，复矩阵秩至多\(2S\)，对应实化矩阵秩通常至多\(4S\)；但多个目标、多个频率和非线性延拓后的并集可以很大。低秩瞬时导数不是全局低秩充分性证明。

这些公式是本轮的任务专用推导，基础伴随微分并非新工具。潜在新颖性在于把它与间隔平滑、完整自散射及背景条件目标分歧联系起来。

\hypertarget{ux4e3aux4ec0ux4e48ux4e0dux80fdux53eaux4fddux7559ux4e00ux9636score}{%
\subsection{11.2
为什么不能只保留一阶score？}\label{ux4e3aux4ec0ux4e48ux4e0dux80fdux53eaux4fddux7559ux4e00ux9636score}}

对\(H(t)=H+tE\)， \[
F_X'=BT_XET_Xa,\qquad
F_X''=2BT_XET_XET_Xa.
\] 所以 \[
(\log L_X)''
=\operatorname{Re}(r_X^*F_X'')-\|F_X'\|^2.
\] 若全路径\(\|T_X\|\le M_T\)、\(\|r_X\|\le R\)，则 \[
|(\log L_X)''|
\le
\left[
2R\|B\|\|a\|M_T^3+
\|B\|^2\|a\|^2M_T^4
\right]\|E\|^2.
\] contrast二阶界还需加anchor对应项。定义 \[
M_D(E)=\sup_{t\in[0,1]}\left|\frac{d^2}{dt^2}D(H+tE)\right|,
\] 其中已包含\(\|E\|^2\)。于是 \[
|D(H+E)-D(H)|
\le|\langle\Xi_H,E\rangle|+\tfrac12M_D(E)
\] 在具有可靠路径稳定界时成立。

当前\(\langle\Xi_H,E\rangle\approx0\)并不意味着有限误差无害。共振附近二阶项会放大；新形状、材料或频率也可能旋转目标电流，激活先前无关的返回方向。

\hypertarget{ux6750ux6599ux76eeux6807ux6df7ux5408ux5bfcux6570}{%
\subsection{11.3
材料--目标混合导数}\label{ux6750ux6599ux76eeux6807ux6df7ux5408ux5bfcux6570}}

对总场固定外源、固定观测算子\(O\)及加性模型\(\mathcal A(W,X)=\mathcal A_0-M_W(W)-M_X(X)\)，其中\(\mathcal A_{WX}=0\)，
\[
F_W=O\mathcal A^{-1}M_Wu,\qquad
F_X=O\mathcal A^{-1}M_Xu,
\] \[
F_{WX}=O\mathcal A^{-1}
(M_W\mathcal A^{-1}M_X+
M_X\mathcal A^{-1}M_W)u.
\] 相应对数likelihood混合导数为 \[
\boxed{
\partial_W\partial_X\log L
=
\operatorname{Re}\langle r,F_{WX}\rangle
-\operatorname{Re}\langle F_W,F_X\rangle.
}
\]
若含交叉几何参数、参数相关观测或非加性材料模型，应增加相应混合导数项。沿\(X_0\)到\(X\)的合法路径积分此式，可以控制墙变化对contrast的影响。两项都不可少：混合传播曲率与一阶响应重合分别描述不同机制。

跨支撑拓扑不能用一个局部路径界全包；应分支处理，保留组间残余质量或统一目标结论证书。

\hypertarget{femux5b9eux73b0ux4e2dux7684ux53f3ux7aefux5fc5ux987bux5b8cux6574}{%
\subsection{11.4
FEM实现中的右端必须完整}\label{femux5b9eux73b0ux4e2dux7684ux53f3ux7aefux5fc5ux987bux5b8cux6574}}

已有后端求散射场： \[
(L_W-V_X)u_{\rm sc}=b_W+b_X.
\]
这里\(b_X\)是目标区域上解析入射场的真实弱式载荷。若\(\mathcal A(z)u(z)=b(z)\)，
\[
u_i=\mathcal A^{-1}(b_i-\mathcal A_i u),
\] \[
u_{ij}=\mathcal A^{-1}
(b_{ij}-\mathcal A_{ij}u-\mathcal A_i u_j-\mathcal A_j u_i).
\]
不能用材料无关右端公式机械代替。只有在合法总场转换后，固定外源假设才成立。

本轮用固定随机复矩阵作了代数核对：\(T\)分离误差约\(7.2\times10^{-16}\)，contrast一阶差分误差约\(1.4\times10^{-12}\)，二阶差分绝对误差约\(1.7\times10^{-7}\)。它只检查公式接线，没有新增PDE证据，也没有认证实际FEM形状导数。

\hypertarget{ux547dux9898dux4ec0ux4e48ux65f6ux5019ux53efux4ee5ux4e0dux8ba1ux7b97ux80ccux666fposteriorux8d28ux91cf}{%
\section{12.
命题D：什么时候可以不计算背景posterior质量？}\label{ux547dux9898dux4ec0ux4e48ux65f6ux5019ux53efux4ee5ux4e0dux8ba1ux7b97ux80ccux666fposteriorux8d28ux91cf}}

\hypertarget{ux4e00ux822cux6982ux7387ux5f15ux7406ux6761ux4ef6ux76eeux6807ux6df7ux5408ux7684ux6536ux7f29}{%
\subsection{12.1
一般概率引理：条件目标混合的收缩}\label{ux4e00ux822cux6982ux7387ux5f15ux7406ux6761ux4ef6ux76eeux6807ux6df7ux5408ux7684ux6536ux7f29}}

将精确目标posterior写为 \[
p_X=\sum_g\beta_gp_g,
\] 其中\(p_g=p(X\mid y,g)\)，\(\beta_g=p(g\mid y)\)。设 \[
\delta=\sup_{g,h}TV(p_g,p_h).
\] 则对任意两套背景权重 \[
\boxed{
TV\left(\sum_g\beta_gp_g,\sum_g\widetilde\beta_gp_g\right)
\le\delta\,TV(\beta,\widetilde\beta)\le\delta.
}
\]

证明：把\(\beta-\widetilde\beta\)的正负部分分别归一化，差即为\(TV(\beta,\widetilde\beta)\)乘两个条件分布凸组合的差；后者距离不超过族直径。这是Dobrushin型一般概率工具，本身不作原创主张。

直接推论：

\begin{itemize}
\tightlist
\item
  条件目标posterior完全相同：背景质量可以完全不算。
\item
  条件posterior接近：允许背景质量较粗。
\item
  条件目标结论明显分歧：相对背景质量必须保留和求准。
\end{itemize}

对连续背景，同样将求和替为积分。

\hypertarget{ux4e0emaxwellux5bf9ux6bd4ux754cux7ec4ux5408}{%
\subsection{12.2
与Maxwell对比界组合}\label{ux4e0emaxwellux5bf9ux6bd4ux754cux7ec4ux5408}}

在共同目标prior下，若可用第10--11节的物理证书证明 \[
\sup_{W,X}|D_W(X)-D_{\rm ref}(X)|\le e_D,
\] 则每个条件posterior与参考之间满足 \[
TV(p_W,p_{\rm ref})\le\tanh(e_D/2).
\]
因此真实目标marginal与\(p_{\rm ref}\)也满足同一界，\textbf{无须知道墙posterior权重}。

这才是可能绕开M4的结构性入口：不是证明未算背景质量很小，而是证明其无论有多少质量，都不会显著改变目标答案。

若墙--目标prior相关，应把\(\log\pi_0(X\mid W)\)的相对变化一起纳入对比。某些背景下anchor不合法时，需要共同支撑分块或多anchor桥接，不能强行使用一个零概率anchor。

\hypertarget{ux51b3ux7b56ux5171ux8bc6ux6bd4posteriorux5171ux8bc6ux66f4ux5f31}{%
\subsection{12.3
决策共识比posterior共识更弱}\label{ux51b3ux7b56ux5171ux8bc6ux6bd4posteriorux5171ux8bc6ux66f4ux5f31}}

若同一动作\(a_0\)在每个条件posterior下的regret都\(\le\epsilon\)，则在任意背景混合权重下，\(a_0\)的regret仍\(\le\epsilon\)。证明为
\[
R_g(a_0)\le \inf_aR_g(a)+\epsilon
\le R_g(a)+\epsilon
\] 后加权并对\(a\)取inf。

平方损失下，条件目标均值相同就足以使Bayes动作相同；方差可以不同。所以该停止规则不能自动输出可靠目标可信区。并且共同动作可能共同不准确：必须同时检验目标绝对误差与信息条件，不能重演G1的``小regret、大RMSE''。

\hypertarget{ux8d28ux91cfux8ba1ux7b97ux7cbeux5ea6ux7684ux6b63ux786eux7075ux654fux5ea6}{%
\subsection{12.4
质量计算精度的正确灵敏度}\label{ux8d28ux91cfux8ba1ux7b97ux7cbeux5ea6ux7684ux6b63ux786eux7075ux654fux5ea6}}

对独立目标prior，先定义精确组边缘likelihood \[
m_g(X)=\int_{\mathcal G_g}L(y\mid W,X)\,\mu_0(dW),\quad
A_g=m_g(X_0),\quad D_g(X)=\log m_g(X)-\log m_g(X_0).
\]
若用代表背景的contrast代替\(D_g\)，下面等式只对近似模型成立，须另加组内包络。写
\[
m(X)=\sum_g A_g e^{D_g(X)},\qquad
q_g(X)=\frac{A_g e^{D_g(X)}}{m(X)}.
\] 则 \[
\boxed{
\frac{\partial}{\partial\log A_g}
[\log m(X)-\log m(X_0)]
=q_g(X)-q_g(X_0).
}
\] 对目标功能\(\phi\)： \[
\frac{\partial}{\partial\log A_g}\mathbb E[\phi(X)\mid y]
=\beta_g\{\mathbb E_g\phi-\mathbb E\phi\}.
\]
上述偏导保持\(D_g\)固定，描述纯组质量误差，不是所有物理墙参数扰动的完整导数。背景质量精度应乘以``不同目标候选/条件功能对该背景的依赖差异''。实际调度必须使用可认证区间上界，不能先读取完整posterior后再计算这个分数。

\hypertarget{ux4e00ux4e2aux5fc5ux987bux4fddux7559ux7684ux53cdux4f8b}{%
\subsection{12.5
一个必须保留的反例}\label{ux4e00ux4e2aux5fc5ux987bux4fddux7559ux7684ux53cdux4f8b}}

两背景、两目标、等prior，likelihood行分别为 \[
L_1=(1,\varepsilon),\qquad L_2=(\varepsilon,1).
\]
目标posterior为\((1/2,1/2)\)。把第二行乘\(\eta\ll1\)，其内部contrast完全不变，但目标posterior趋向
\[
\left(\frac1{1+\varepsilon},\frac{\varepsilon}{1+\varepsilon}\right).
\]
因此逐背景任意删除likelihood尺度通常错误。只有先证明条件目标共识、共同动作，或保留相对质量，才可安全省略。这是概率反例，不冒充已运行Maxwell物理世界。

\hypertarget{ux600eux6837ux5728ux4e0dux904dux5386ux5168ux90e8ux539fux5b50ux65f6ux5f97ux5230ux4fbfux5b9cux8bc1ux4e66}{%
\section{13.
怎样在不遍历全部原子时得到便宜证书？}\label{ux600eux6837ux5728ux4e0dux904dux5386ux5168ux90e8ux539fux5b50ux65f6ux5f97ux5230ux4fbfux5b9cux8bc1ux4e66}}

\hypertarget{ux8bc1ux4e66ux5e94ux6709ux5206ux5c42ux5165ux53e3}{%
\subsection{13.1
证书应有分层入口}\label{ux8bc1ux4e66ux5e94ux6709ux5206ux5c42ux5165ux53e3}}

\begin{enumerate}
\def\labelenumi{\arabic{enumi}.}
\tightlist
\item
  微结构类：用共同局部均值、材料幅度与稳定裕度，先给整个类的传播包络。
\item
  返回算子：用分离间隔、端口尾部或独立action验证控制\(H-\widetilde H\)。
\item
  目标对比：用成对前向--伴随余切和二阶余项，把传播误差变成\(D\)误差。
\item
  概率/决策：用conditional consensus、边缘likelihood区间或风险包络停止。
\item
  如果证书失效，增加mode、拆组或回退完整求解；每次回退计费。
\end{enumerate}

没有物理组结构的任意黑箱likelihood，不可能普遍保证少查询：未算而有非零prior的一个原子，在另一张与已知记录一致的表中可以承载大量后验质量。M4的失败不只是实现粗糙，也反映默认上界缺乏信息。

\hypertarget{ux4e00ux4e2aux6709ux9650ux7ef4ux6982ux7387actionux8bc1ux4e66}{%
\subsection{13.2
一个有限维概率action证书}\label{ux4e00ux4e2aux6709ux9650ux7ef4ux6982ux7387actionux8bc1ux4e66}}

对固定实化离散返回误差矩阵\(E\)，取与构造\(E\)独立的\(g_i\sim N(0,I)\)，\(i=1,\dots,q\)。令
\[
t_{\alpha,q}=\Phi^{-1}\left(\frac{1+\alpha^{1/q}}2\right).
\] 则 \[
\boxed{
\Pr\left\{\|E\|_2
\le\frac{\max_i\|Eg_i\|_2}{t_{\alpha,q}}\right\}\ge1-\alpha.
}
\]
证明：顶右奇向量\(v\)满足\(\|Eg_i\|\ge\|E\||v^Tg_i|\)，而所有\(|v^Tg_i|<t\)的概率为\([2\Phi(t)-1]^q\)。

它可以只用若干真实返回actions验证一个大离散算子，不需要全部2257列。但以下费用和限制不可忽略：

\begin{itemize}
\tightlist
\item
  背景FE分解以及每个wall/frequency的\(q\)个RHS仍收费；
\item
  probe必须覆盖所声称的完整离散输入空间，不能只在训练子空间内采样；
\item
  训练基与验证probe必须独立；适应轮次使用fresh
  probes或同时界，累计\(\sum\alpha_k\le\alpha_{\rm total}\)；
\item
  使用物理迹范数时要作质量矩阵加权；
\item
  数值action误差需加到\(\|Eg_i\|\)上界；求解残差没有稳定界时不是场误差证书；
\item
  该概率只针对验证随机性，不是Bayes可信度或实验覆盖率；
\item
  它只认证声明离散空间，不能覆盖遗漏的连续模式、PML和离散误差；
\item
  高秩时界可能很保守。正确的证书也可能没有经济价值。
\end{itemize}

\hypertarget{ux4eceux5bf9ux6bd4ux8befux5deeux5230ux76eeux6807posterior}{%
\subsection{13.3
从对比误差到目标posterior}\label{ux4eceux5bf9ux6bd4ux8befux5deeux5230ux76eeux6807posterior}}

采用同一个目标prior、在共同支持上比较时，一般概率界仍为 \[
\operatorname{osc}_X\log\frac{\widetilde m_y(X)}{m_y(X)}
\le\kappa
\quad\Longrightarrow\quad
TV(\widetilde p_X,p_X)\le\tanh(\kappa/4).
\] 它是已有工具，不列为本轮独立Maxwell定理。

若损失\(0\le\ell\le L_{\max}\)，TV为\(t\)时，任一固定动作的风险误差不超过\(L_{\max}t\)，近似posterior的Bayes动作在精确posterior下的regret不超过\(2L_{\max}t\)。若是直径为\(D_{\rm ROI}\)的紧目标位置域上的平方损失，采用各自posterior均值，则更有
\[
\operatorname{Regret}\le D_{\rm ROI}^2t^2.
\] 这些是相对于精确posterior的计算风险界，不是对真实目标误差的保证。

若TV门为0.02，则\(\kappa\)约需不超过0.08001。直接的0.1前向白化误差并不自动足够：Gaussian
potential还乘以残差，必须显式转换。

对有限目标原子，若未归一质量有\(a_j\le w_j\le b_j\)，事件\(E\)满足 \[
\frac{\sum_Ea_j}{\sum_Ea_j+\sum_{E^c}b_j}
\le P(E\mid y)\le
\frac{\sum_Eb_j}{\sum_Eb_j+\sum_{E^c}a_j}.
\]
区间分母须为正，\(\sum_j a_j>0\)是充分条件；若某分母为零，应取保守下界0或上界1，不能直接相除。背景残余质量、目标分支遗漏和数值误差必须进入同一个包络。ESS或已见样本稳定不能替代未知质量上界。若传播误差来自随机probe，后续posterior与决策界只在数值验证成功事件上成立，整体成功概率至少为\(1-\alpha_{\rm total}\)，不能称为无条件确定性证书。

\hypertarget{ux4e09ux6761ux4e0dux540cux7684ux7b97ux6cd5ux8defux7ebf}{%
\section{14.
三条不同的算法路线}\label{ux4e09ux6761ux4e0dux540cux7684ux7b97ux6cd5ux8defux7ebf}}

\hypertarget{ux8defux7ebfiux5206ux8fa8ux7387ux5339ux914dux7684ux6709ux6548ux4f20ux64adux53cdux6f14}{%
\subsection{路线I：分辨率匹配的有效传播反演}\label{ux8defux7ebfiux5206ux8fa8ux7387ux5339ux914dux7684ux6709ux6548ux4f20ux64adux53cdux6f14}}

\textbf{解决信息与表示问题。}
在普通二维微结构合同下，不尝试逐格恢复细墙。用未知宏观传播参数\(\zeta\)、完整目标参数\(X\)与共享\(\gamma\)联合拟合或积分，并把微结构误差作为有界/有来源的模型不确定性。

步骤：冻结宏观类与微结构包络；验证有效全波误差；用measurement-only初值进行充分的多起点完整非线性反演；保留导致不同目标的宏观分支；输出目标估计、有限补偿证据和不确定性。

它可能通过降低病态的材料自由度改善实际求解；这不是数据凭空增加。若宏观类漏掉大空腔或形状变化，应判为模型不充分，而不是强迫目标吸收误差。

关键对照：全分辨率联合FWI、普通有效介质FWI、已知真实墙oracle、错误固定墙、带同样模型误差处理的基线。

\hypertarget{ux8defux7ebfiiux72ecux7acbux76eeux6807ux6563ux5c04ux5899ux8fd4ux56deux7ec4ux5408}{%
\subsection{路线II：独立目标散射＋墙返回组合}\label{ux8defux7ebfiiux72ecux7acbux76eeux6807ux6563ux5c04ux5899ux8fd4ux56deux7ec4ux5408}}

\textbf{解决物理求解复用问题。}
预备墙的边界返回与入射/接收actions，独立处理目标自散射模板。墙--目标组合只解小闭环，按目标contrast证书适应增加返回mode。

它与普通两物体T-matrix、Schur、optimal port
reduction同源。新增比较必须是：同一返回表示下，目标概率驱动精度相对统一输出精度是否省真实工作。路线II对完整端口的优势不等于Bayesian部分的原创收益。

适合固定圆盘平移等有可复用内禀模板的合同；目标靠墙、强共振、复杂新形状可使秩和模板成本增加。

\hypertarget{ux8defux7ebfiiiux6761ux4ef6ux76eeux6807ux5171ux8bc6ux9a71ux52a8ux7684ux8fb9ux7f18ux63a8ux65ad}{%
\subsection{路线III：条件目标共识驱动的边缘推断}\label{ux8defux7ebfiiiux6761ux4ef6ux76eeux6807ux5171ux8bc6ux9a71ux52a8ux7684ux8fb9ux7f18ux63a8ux65ad}}

\textbf{解决概率质量计算问题。}
先用物理包络组织背景组，求代表的条件目标对比及组内误差。然后：

\begin{itemize}
\tightlist
\item
  全部条件目标posterior接近：直接给有界目标marginal近似；
\item
  仅共同动作一致：只交付决策与regret证书；
\item
  条件目标分歧：求这些组的相对质量，优先收紧最可能改变目标结论的组；
\item
  无法取得便宜物理包络：退回普通自适应求积，不声称节省。
\end{itemize}

这条路线与普通goal-ROM的区别不是少几个输出，而是\textbf{允许背景响应和背景质量很不准确，只要它们对归一化目标结论无效}。但若每个组员都要先完整算anchor才能验证共识，就只省了后续精度，不能把anchor
FE算成已避免。

\hypertarget{ux5907ux7528ux8defux7ebfux6570ux636eredatumingux4e0eux8c31ux4efbux52a1}{%
\subsection{备用路线：数据redatuming与谱任务}\label{ux5907ux7528ux8defux7ebfux6570ux636eredatumingux4e0eux8c31ux4efbux52a1}}

富宽带、充分采样、已知direct-arrival/macromodel的合同，可研究Marchenko或数据驱动EM
ROM直接获得目标边界场，再作局部反演。{[}R3--R6,R9{]}
这是不同的数据路线，当前单频bank不足以验证。

另一备用是目标谱功能：孤立目标共振在有界返回扰动下的极点漂移。但目标自然共振通常依赖背景；已有埋藏介质目标TAP结果展示深度变化导致极点轨迹。{[}R18{]}
它首先可能服务尺寸/材料组合，不服务自由空间绝对位置。详见附录A，不作为下一轮默认PDE任务。

\hypertarget{ux5b9eux9645ux6210ux672cux4ec0ux4e48ux7ed3ux6784ux53efux80fdux6bd4ux5b8cux6574greenux66f4ux4fbfux5b9c}{%
\section{15.
实际成本：什么结构可能比完整Green更便宜？}\label{ux5b9eux9645ux6210ux672cux4ec0ux4e48ux7ed3ux6784ux53efux80fdux6bd4ux5b8cux6574greenux66f4ux4fbfux5b9c}}

令\(H\)个墙、\(J\)个目标、每频\(S\)个源。完整逐对求解近似成本为 \[
C_{\rm full}=HJ(C_{\rm FE}+S C_{\rm RHS}).
\] 分离返回表示的一种成本模型为 \[
\begin{aligned}
C_{\rm split}={}&
H(C_W^{\rm FE}+r C_W^{\rm RHS}+S C_W^{\rm RHS})\\
&+K C_{\rm target\ template}
+HJ C_{\rm small}(r)\\
&+C_{\rm certificate}+C_{\rm fallback}.
\end{aligned}
\]
\(K\)为内禀目标模板数，通常不等于位置数；固定圆盘每频可\(K=1\)。若每个目标都需\(r\)个模板RHS，则总RHS约\((H+J)r+HS\)。

旧\(H=118,J=25,S=2\)时，逐对前向为5900
RHS；简化分离账为\(143r+236\)，仅在RHS上占优要求\(r<39.6\)，还未加证书、伴随和回退。这是代数break-even，不是对实际秩的预测。固定圆盘模板复用会改变该账，但所有对照必须同享。

\textbf{旧合同两个源RHS已经一次给出全部16接收输出。}
把32个复响应压成一个contrast标量，本身不减少FE/RHS，增加伴随还可能更贵。

报告必须分别给出：

\begin{itemize}
\tightlist
\item
  logical query；
\item
  实际complex RHS；
\item
  材料--频率FE分解；
\item
  small dense分解；
\item
  setup、证书与fallback CPU；
\item
  wall time、峰值RSS和常驻/磁盘缓存；
\item
  冷启动、有限重复任务摊销两种合同。
\end{itemize}

若只有普通分离T-matrix获益，而概率调度无增量，应据实归因。若只减少查表、不减少部署物理费用，应标为逻辑查询收益。

\hypertarget{ux4e09ux7ec4ux51b3ux5b9aux6027ux5b9eux9a8c}{%
\section{16.
三组决定性实验}\label{ux4e09ux7ec4ux51b3ux5b9aux6027ux5b9eux9a8c}}

\hypertarget{e1wrong-wall-correct-target}{%
\subsection{E1：Wrong Wall, Correct
Target}\label{e1wrong-wall-correct-target}}

\hypertarget{e1-pux7269ux7406ux5b58ux5728ux6027}{%
\subsubsection{E1-P：物理存在性}\label{e1-pux7269ux7406ux5b58ux5728ux6027}}

先冻结三类互补世界。

\begin{enumerate}
\def\labelenumi{\arabic{enumi}.}
\tightlist
\item
  \textbf{二维微结构主臂：}
  相同局部复介电均值，不同二维纹理或小孔隙；至少三个可被独立网格解析的尺度。材料L2差保持明显，完整目标族响应差随尺度下降。
\item
  \textbf{镜像结构正对照：}
  二维非均匀对称墙、对称空腔、成对端口；验证奇偶信息分离及非对称破坏。
\item
  \textbf{有限补偿反例：}
  在严格被动墙内选真实产品核方向，检验满射并沿目标变化作全波正确器；固定新增端口/频率检查其失效。
\end{enumerate}

目标族至少包含位置、半径、实部、损耗变化及圆/椭圆两形状。先用少量anchor构造，再以事前冻结而未用于构造的目标检验。不能为每个新目标重新调墙后称同一传播等价。

真值FE必须解析微结构；粗网格数值平均不是物理等价证据。至少一个强目标要求Born与full响应显著不同，以确认完整反馈确实被检验。

\hypertarget{e1-oux5728ux7ebfux672aux77e5ux76eeux6807}{%
\subsubsection{E1-O：在线未知目标}\label{e1-oux5728ux7ebfux672aux77e5ux76eeux6807}}

算法只见\(y\)、公开ROI、材料类与proper
prior。真实平均墙、实际mask、目标位置及oracle支撑均不输入。

同时比较有效传播推断、充分多起点联合FWI、普通两阶段墙估计、真实墙oracle。训练、在线输出、真值评分分开。固定预算触限要标为优化未达，不解释为物理不存在。

\hypertarget{e1-gux4efbux52a1ux6269ux5c55ux4e0eux7834ux574f}{%
\subsubsection{E1-G：任务扩展与破坏}\label{e1-gux4efbux52a1ux6269ux5c55ux4e0eux7834ux574f}}

测试更大目标、不同形状、更高频、更近墙位置、低损耗强闭环及宏观空腔。分别判断来自微结构不再小、间隔缩小、闭环放大、宏观模型缺项、仪器gauge或优化失败。

\textbf{新主定位门建议事前冻结为RMSE\(\le0.1\lambda_0\)；仅为下一轮研究门，不是当前实测。}
同时报告90分位与最坏误差，并与zero-data
prior和oracle比较。尺寸、介电率、损耗及形状各自单列，不用位置成功替代完整目标成功。

\hypertarget{e2ux76eeux6807ux53efux8fa8ux6027ux4e0eux5899ux53efux8fa8ux6027}{%
\subsection{E2：目标可辨性与墙可辨性}\label{e2ux76eeux6807ux53efux8fa8ux6027ux4e0eux5899ux53efux8fa8ux6027}}

对E1同一批世界，严格嵌套使用：

\begin{itemize}
\tightlist
\item
  \(\mathcal A_0\)：原两发16收、\(f_0\)；
\item
  \(\mathcal A_1\)：保留原数据，加入两个外侧源、同16收、\(f_0\)；
\item
  \(\mathcal A_2\)：保留以上数据，加入\(0.85f_0,1.15f_0\)的四源16收；
\item
  更高频、非对称仪器等作为独立stress合同。
\end{itemize}

每例同时报告\(K_{W\mid X,\gamma}\)、\(K_{X\mid W,\gamma}\)、绝对奇异值、数值不确定性、有限幅度补偿曲线、非线性\(\omega_X\)见证下界，以及可认证区域内的上界。形状分支用有限secant。

目标参数需按声明单位尺度化。精确秩、数值阈值秩、约束有效维数分列。核查新增数据后目标补偿集是否相交，而不是只比较相对谱比例。

真实目标MSE、posterior风险、事件概率、形状正确率、覆盖与预测评分是主指标；wall方差和L2误差只作诊断。至少寻找一类``目标风险显著降、wall细节仍宽''的完整全波世界。

\hypertarget{e3ux4f3cux7136ux5bf9ux6bd4ux4e0egreen-actionux6210ux672c}{%
\subsection{E3：似然对比与Green-action成本}\label{e3ux4f3cux7136ux5bf9ux6bd4ux4e0egreen-actionux6210ux672c}}

\hypertarget{ux7b2cux4e00ux9636ux6bb5ux65e7bankux53d7ux63a7ux903bux8f91ux5f00ux53d1}{%
\subsubsection{第一阶段：旧bank受控逻辑开发}\label{ux7b2cux4e00ux9636ux6bb5ux65e7bankux53d7ux63a7ux903bux8f91ux5f00ux53d1}}

只开放prior、\(y\)和按需主端口响应oracle。完整响应表、完整posterior、truth、外侧TX2/3评分数据不能被在线进程预读。所有策略输出封存后再生成reference与评分。

比较四种组距离：完整响应、目标增量、含残差投影的contrast、条件目标posterior/决策。连续目标推广不从25个位置自动得出。

至少实现：

\begin{itemize}
\tightlist
\item
  完整2950枚举；
\item
  普通自适应求积/默认质量包络；
\item
  普通goal-ROM或同等response插值；
\item
  普通分离返回表示、统一输出门；
\item
  同表示下contrast精度与条件目标共识策略。
\end{itemize}

旧bank不含新返回actions或算子导数；需要新物理的证书不能凭全表离线拟合冒充在线已知。任何teacher/offline统计都在成本单列并给所有基线同权。

\hypertarget{ux7b2cux4e8cux9636ux6bb5ux6709ux9650ux51b7ux542fux52a8pde}{%
\subsubsection{第二阶段：有限冷启动PDE}\label{ux7b2cux4e8cux9636ux6bb5ux6709ux9650ux51b7ux542fux52a8pde}}

只有逻辑实验显示非平凡查询减少、质量门通过且证书可获得，才做小规模冷启动。建议8个新墙、25个冻结目标原子，包含至少两种内禀模板，逐步取\(r=4,8,16,32\)，按质量停止或fallback。使用新观测与独立场资格，声明这是有限prior验证。

比较实际FE/RHS/CPU和峰值内存。完整baseline与新法每次从相同空缓存起步；共同合法物理模板对所有方法开放。评估reference成本与方法部署成本分开记，实际总花费也完整记录。

\hypertarget{ux6982ux7387ux4e0eux51b3ux7b56ux95e8}{%
\subsubsection{概率与决策门}\label{ux6982ux7387ux4e0eux51b3ux7b56ux95e8}}

\begin{itemize}
\tightlist
\item
  目标TV、事件/CDF误差建议\(\le0.02\)；
\item
  所有置信界越界与数值失败保留；
\item
  真值定位门独立于对reference的regret；
\item
  新prior预测重复试验给风险和覆盖区间；少量parent只作pilot，不称正式SBC；
\item
  证书随机性、观测噪声和prior重复是三个不同随机源；
\item
  若无净成本收益，明确失败来自秩、证书保守、质量不一致或setup，不改变质量门。
\end{itemize}

\hypertarget{ux539fux521bux6027ux4e0eux6700ux5f3aux8fd1ux90bb}{%
\section{17.
原创性与最强近邻}\label{ux539fux521bux6027ux4e0eux6700ux5f3aux8fd1ux90bb}}

\begin{longtable}[]{@{}
  >{\raggedright\arraybackslash}p{(\columnwidth - 4\tabcolsep) * \real{0.3333}}
  >{\raggedright\arraybackslash}p{(\columnwidth - 4\tabcolsep) * \real{0.3333}}
  >{\raggedright\arraybackslash}p{(\columnwidth - 4\tabcolsep) * \real{0.3333}}@{}}
\toprule\noalign{}
\begin{minipage}[b]{\linewidth}\raggedright
近邻
\end{minipage} & \begin{minipage}[b]{\linewidth}\raggedright
已经覆盖
\end{minipage} & \begin{minipage}[b]{\linewidth}\raggedright
本研究还需交付
\end{minipage} \\
\midrule\noalign{}
\endhead
\bottomrule\noalign{}
\endlastfoot
DBIM，Chew--Wang 1990 {[}R1{]} & 更新背景处理超出Born的材料反演 &
目标nuisance下的有限辨识与概率充分性 \\
Watson等2024 {[}R7{]} &
墙校正Green、未知墙参数、ROM，错误墙几何下改善目标图像 &
完整目标反馈、目标稳定定理和可靠marginal \\
Cakoni等2019 {[}R2{]} & 二维零阶微结构均匀化与边界修正 &
对未知目标族统一、反馈与目标分辨率联结 \\
Marchenko target replacement / FWI {[}R3,R4{]} &
在充分波场采样及direct-arrival/平滑宏观模型合同下，不需详细外域材料即可处理完整多次反射和目标局部更新
& 本稀疏有限孔径Maxwell合同的信息与概率条件 \\
耗散EM Marchenko {[}R5,R6{]} &
耗散EM内部Green获取已有理论；所引结果依赖双侧反射/透射信息，所引耗散数值验证为层状或一维
& 单侧二维非均匀未知目标的可靠性 \\
Smetana--Patera、Haynes--Fenni {[}R8,R10{]} & 最优端口与多体T-matrix压缩
& 目标contrast使哪些计算确实可省 \\
Borcea等2024 {[}R9{]} & 无介质真值的数据驱动EM ROM及内部波场 &
固定物理坐标gauge与目标nuisance posterior \\
Manzoni等2016 {[}R11{]} & 不可辨nuisance边缘化、RB与后验误差模型 &
强反馈Maxwell专用便宜证书及FE/RHS收益 \\
Spantini等2017 {[}R12{]} & 直接目标posterior、线性Gaussian最优降阶 &
非Gaussian多分支与非线性全波 \\
Zahm等2022 {[}R13{]} & 非线性likelihood降维及KL认证 &
对比消去共同尺度、证书实际经济性 \\
Kaipio等2019 {[}R14{]} & 用Bayesian approximation error改进Born &
精确nuisance边缘化与固定世界可靠性 \\
Porat--Vinnikov {[}R15{]} & 非交换有理函数最小实现 &
不能将静态物理words当独立可激励探针 \\
自适应Bayesian求积 {[}R16{]} & 主动选节点与收敛理论 &
Maxwell物理组包络优于通用插值的证据 \\
\end{longtable}

还检索到Wang等2026预印本{[}R17{]}，摘要涉及分离支撑、有限模态、collective
resolvent与算子范数收敛；本轮核对到摘要，未核全文证明。因此不对其所有定理作排他判断，但它足以说明``分离支撑＋返回模态收敛''必须认真比对。

上述调查不是全球新颖性认证。可以争取的原创组合是：

\[
\text{普通二维可验证传播机制}
\longrightarrow
\text{目标有限补偿/contrast界}
\longrightarrow
\text{不需要的背景质量或actions确实被跳过}.
\]

\hypertarget{ux5206ux5408ux540cux7684go-no-goux88c1ux51b3}{%
\section{18. 分合同的GO /
NO-GO裁决}\label{ux5206ux5408ux540cux7684go-no-goux88c1ux51b3}}

\begin{longtable}[]{@{}
  >{\raggedright\arraybackslash}p{(\columnwidth - 4\tabcolsep) * \real{0.3333}}
  >{\raggedright\arraybackslash}p{(\columnwidth - 4\tabcolsep) * \real{0.3333}}
  >{\raggedright\arraybackslash}p{(\columnwidth - 4\tabcolsep) * \real{0.3333}}@{}}
\toprule\noalign{}
\begin{minipage}[b]{\linewidth}\raggedright
具体主张
\end{minipage} & \begin{minipage}[b]{\linewidth}\raggedright
裁决
\end{minipage} & \begin{minipage}[b]{\linewidth}\raggedright
下一项决定性证据
\end{minipage} \\
\midrule\noalign{}
\endhead
\bottomrule\noalign{}
\endlastfoot
精确共同离散背景消元 & 已成立 & 复用后端，不重复当主科研 \\
普通二维细微结构材料差大、宽任务响应近 & 条件理论GO &
微结构解析全波及目标族统一性 \\
任意墙类、有限数据总能恢复目标 & NO-GO & 全波产品核满射给反例 \\
对称墙下特定横向任务局部可辨 & 条件理论GO & 奇偶分解与破坏对照 \\
精确Schur为零就没有实用目标分辨率 & 该推断错误 &
有限幅度、噪声尺度profile \\
整个目标Green普遍低秩 & NO-GO作为普遍主张 & 丰富目标电流可满维 \\
只压缩分离后的墙返回 & 理论GO，成本未证 &
实际模式尾、闭环放大及冷账本 \\
目标增量相同即可删背景权重 & NO-GO & 残差投影与条件目标分歧反例 \\
条件目标共识可省质量精度 & 条件数学GO & 便宜Maxwell共识证书 \\
当前数据contrast相同即宽任务等效 & NO-GO & 噪声、目标、孔径扩展检验 \\
当前BGGA已实现目标准确且加速 & NO-GO &
绝对误差和实际避免FE/RHS均未成立 \\
新的似然分配必然胜普通ROM & 待验证 & 相同缓存同质量的净收益 \\
A理论主线可继续，即使C不加速 & GO & 物理定理与决定性正反全波证据 \\
\end{longtable}

负面结果必须归因到正确层次：信息论歧义、模型类错误、观测不足、gauge过宽、表示不合适、后验不能压缩、证书太贵、非凸求解未达，不能互相替代。

\hypertarget{ux6700ux503cux5f97ux5f62ux6210ux5f3aieee-tapux8bbaux6587ux7684ux5355ux4e00ux8d21ux732e}{%
\section{19. 最值得形成强IEEE
TAP论文的单一贡献}\label{ux6700ux503cux5f97ux5f62ux6210ux5f3aieee-tapux8bbaux6587ux7684ux5355ux4e00ux8d21ux732e}}

建议工作题目：

\textbf{Target Resolution through Unresolved Dielectric Walls: Nonlinear
Maxwell Propagation Classes and Task-Stability Bounds}

中文：\textbf{不可分辨介电墙体后的目标分辨率：非线性Maxwell传播类与任务稳定界。}

主结果应连成一条线，而不是堆叠八种工具。

\begin{enumerate}
\def\labelenumi{\arabic{enumi}.}
\tightlist
\item
  \textbf{物理正机制：}
  普通二维墙材料差可保持大，而声明目标族传播差受微结构尺度、间隔与闭环稳定控制。
\item
  \textbf{物理边界：}
  合法全波墙产品核满射时存在精确目标补偿；对称性或可验证有限宏观分离可保护特定任务。
\item
  \textbf{有限目标界：}
  将真实可算的传播误差与目标分离裕度联结，得到分辨率和新增测量的作用范围。
\item
  \textbf{实际算法：}
  在不知道真值墙和目标时，联合推断有效传播与目标；若条件目标结论分歧则保留概率分支。
\item
  \textbf{决定性实验：}
  绝对目标误差低、墙细节仍不可恢复；反例在增加非对称自由度、接近共振或扩大任务时按理论出现。
\end{enumerate}

若E3另外证明真实FE/RHS/CPU净减少，可把contrast计算作为强算法增量；若没有，不能用查表速度或放宽门槛救算法叙事。A的物理问题仍值得独立完成。

仅有均匀化＋假设逆Lipschitz，可能不足以支撑强TAP。至少还需一个有具体几何解释、经过完整全波验证的目标分离机制，或一个在普通宏观墙不确定性中仍有效的有限补偿界。E1镜像正对照用于验证机制，不能代替更宽的目标任务实验。

\hypertarget{ux5bf9ux6700ux7ec8ux79d1ux5b66ux95eeux9898ux7684ux56deux7b54}{%
\section{20.
对最终科学问题的回答}\label{ux5bf9ux6700ux7ec8ux79d1ux5b66ux95eeux9898ux7684ux56deux7b54}}

\textbf{当墙体不可辨识时，需要的不是准确墙体材料图，而是足以区分合法目标解释的传播作用。}

物理层面，这通常包括目标域有效照明、到接收器的可观测传播，以及会返回目标并被其再次散射的反馈；只在目标可达、可观测的方向上需要精确。目标自散射可以完整保留，墙返回可以在有条件的误差预算内压缩。

概率层面，必须保留会改变目标相对likelihood的背景变化。对所有条件目标结论共同无效的背景尺度和质量可以省；造成不同目标位置、材料或形状解释的分支必须积分或保留。单Green不一般闭合。

信息层面，目标可以比墙细节更易辨识，但并不总比墙的某些宏观传播参数更易辨识。新增测量只需打破目标所需的共同墙补偿，不必恢复墙全部材料。无限小精确歧义、有限噪声分辨率和proper-prior概率风险应同时报告。

计算层面，最有实质希望的是：\textbf{目标自散射模板复用＋平滑墙返回actions＋目标contrast/条件共识决定精度和质量计算。}
其中前两项已有成熟物理方法；本研究的新价值必须由后者的可验证收益和前者的目标稳定解释来体现。

当前证据仍未证明一般二维未知墙下的可靠目标成像与净加速已经实现。本轮给出了普通二维正机制、真实非线性反例、明确误差链和可执行检验，下一步应把它们变成E1/E2/E3的决定性数据。

\hypertarget{ux9644ux5f55aux5171ux632fux4efbux52a1ux7684ux6761ux4ef6ux6027ux5907ux7528ux7ed3ux8bba}{%
\section{附录A：共振任务的条件性备用结论}\label{ux9644ux5f55aux5171ux632fux4efbux52a1ux7684ux6761ux4ef6ux6027ux5907ux7528ux7ed3ux8bba}}

设孤立目标\(K_0(\omega)=I-C_0(\omega)V_X(\omega)\)在\(\omega_0\)有简单零点，左右向量\(l,r\)满足\(l^*K_0'(\omega_0)r\ne0\)。墙返回扰动为\(E_W=H_WV_X\)，则
\[
\omega_t-\omega_0
=t\frac{l^*E_W(\omega_0)r}{l^*K_0'(\omega_0)r}+O(t^2).
\] 这是已有解析特征值微扰工具，不作新颖性主张。

保证域内零点总重数保持不变的一个可验证充分条件是：在相关算子于闭域内解析、轮廓上可逆并满足Fredholm条件时，
\[
\sup_{\omega\in\partial D}
\|K_0(\omega)^{-1}H_W(\omega)V_X(\omega)\|<1,
\] 这不是必要条件。它先控制域内零点总重数；小位移还需隔离半径与条件数。

观测留数 \[
\operatorname{Res}_{\omega_p}F
=
\frac{B_WV_X r\,l^*a_W}{l^*K'(\omega_p)r}
\]
必须非零且足够大，墙传播在附近无额外极点。实频弱返回或大间距不能代替复频轮廓控制；开放共振态在下半平面可能随空间距离指数增长。{[}R18,R19{]}

自由空间目标极点不随平移改变，所以不能单独恢复目标位置。有限噪声的谱峰、拟合极点与真实复频极点不同。这条路线仅保留为目标尺寸/材料/形状功能的研究备选。

\hypertarget{ux53c2ux8003ux6587ux732eux4e0eux56faux5b9aux6765ux6e90}{%
\section{参考文献与固定来源}\label{ux53c2ux8003ux6587ux732eux4e0eux56faux5b9aux6765ux6e90}}

\hypertarget{ux56faux5b9aux7814ux7a76ux8bc1ux636e}{%
\subsection{固定研究证据}\label{ux56faux5b9aux7814ux7a76ux8bc1ux636e}}

{[}E1{]} Research W-BGGA，73ba6eec。
\href{https://github.com/migodam/b1-probabilistic-maxwell-imaging/blob/73ba6eec3746548bf21259f627b635678709df00/research/research_w_bgga/RESEARCH_W_BGGA_REPORT_ZH.md}{固定报告}；
\href{https://github.com/migodam/b1-probabilistic-maxwell-imaging/blob/73ba6eec3746548bf21259f627b635678709df00/research/research_w_bgga/GATE_DECISION.json}{机器判决}；
\href{https://github.com/migodam/b1-probabilistic-maxwell-imaging/blob/73ba6eec3746548bf21259f627b635678709df00/research/research_w_bgga/COST_SUMMARY.json}{最终成本}。

{[}E2{]} Research W-BGGA，6fd50478。
\href{https://github.com/migodam/b1-probabilistic-maxwell-imaging/blob/6fd504785f4e4be74d919f2363e0b2d1ef111534/research/research_w_bgga/RESEARCH_W_BGGA_REPORT_ZH.md}{固定第一版报告}。

{[}E3{]} Research W，958a4787。
\href{https://github.com/migodam/b1-probabilistic-maxwell-imaging/blob/958a478739d9553269316ded8125fda9df225c0c/research/research_w_p0r_to_p1/RESEARCH_REPORT_ZH.md}{固定报告}。

{[}E4{]} Research W-SLAM，52351aee。
\href{https://github.com/migodam/b1-probabilistic-maxwell-imaging/blob/52351aee6ddd41b09851fe5e52d1a8861891ef13/research/research_w_slam/RESEARCH_W_SLAM_REPORT_ZH.md}{固定报告}。

{[}E5{]} Research
W-GQ深度研究报告，2026-10-28，35页。已读取完整正文、参考文献和Codex附录；属于上一轮理论交付，不是新增实验。

\hypertarget{ux6700ux5f3aux8fd1ux90bbux53caux5de5ux5177ux6765ux6e90}{%
\subsection{最强近邻及工具来源}\label{ux6700ux5f3aux8fd1ux90bbux53caux5de5ux5177ux6765ux6e90}}

{[}R1{]} Chew W. C.; Wang Y. M. (1990). Reconstruction of
two-dimensional permittivity distribution using the distorted Born
iterative method. IEEE TMI 9(2):218--225.
\href{https://hub.hku.hk/handle/10722/182506}{DOI与作者机构记录}，10.1109/42.56334。

{[}R2{]} Cakoni F.; Guzina B. B.; Moskow S.; Pangburn T. (2019).
Scattering by a Bounded Highly Oscillating Periodic Medium and the
Effect of Boundary Correctors. SIAM JAM 79(4):1448--1474.
\href{https://epubs.siam.org/doi/10.1137/19M1237089}{原文页}。

{[}R3{]} Wapenaar K.; Staring M. (2018). Marchenko-Based Target
Replacement, Accounting for All Orders of Multiple Reflections. JGR
123:4942--4964.
\href{https://doi.org/10.1029/2017JB015208}{DOI}，\href{https://arxiv.org/abs/1707.06580}{作者预印本}。

{[}R4{]} Cui T.; Rickett J.; Vasconcelos I.; Veitch B. (2020).
Target-oriented full-waveform inversion using Marchenko redatumed
wavefields. GJI 223(2):792--810.
\href{https://research-portal.uu.nl/ws/portalfiles/portal/235427686/ggaa333.pdf}{作者机构全文}，10.1093/gji/ggaa333。

{[}R5{]} Slob E.; Zhang L. (2016). Electromagnetic Marchenko equations
for a dissipative heterogeneous medium. ICGPR.
\href{https://research.tudelft.nl/en/publications/electromagnetic-marchenko-equations-for-a-dissipative-heterogeneo/}{作者机构摘要}，10.1109/ICGPR.2016.7572702。

{[}R6{]} Yang B.; Slob E. (2021). Marchenko Inversion of GPR Data for a
1D Dissipative Medium. PIER M 102:65--79.
\href{https://doi.org/10.2528/PIERM21020901}{DOI}。

{[}R7{]} Watson F. M.; Andre D.; Lionheart W. R. B. (2024). Resolving
Full-Wave Through-Wall Transmission Effects in Multi-Static Synthetic
Aperture Radar. Inverse Problems.
\href{https://arxiv.org/html/2403.10354v1}{原文}，10.1088/1361-6420/ad5b83。

{[}R8{]} Haynes M. S.; Fenni I. (2022). Compressed T-Matrix Algorithm
for Scalar and Electromagnetic Scattering from Multiple Objects and
Multiple Incident Directions. PIER M 110:61--72.
\href{https://www.jpier.org/issues/volume.html?paper=22030901}{期刊入口}，10.2528/PIERM22030901。

{[}R9{]} Borcea L.; Liu Y.; Zimmerling J. (2024). Electromagnetic
inverse wave scattering in anisotropic media via reduced order modeling.
JCP 515:113272.
\href{https://arxiv.org/html/2403.03844v1}{作者全文}，10.1016/j.jcp.2024.113272。

{[}R10{]} Smetana K.; Patera A. T. (2016). Optimal Local Approximation
Spaces for Component-Based Static Condensation Procedures. SISC
38(5):A3318--A3356.
\href{https://dspace.mit.edu/entities/publication/98c5582d-4c97-4e6e-9f02-2bea41a29a29}{作者机构原文}，10.1137/15M1009603。

{[}R11{]} Manzoni A.; Pagani S.; Lassila T. (2016). Accurate Solution of
Bayesian Inverse Uncertainty Quantification Problems Combining Reduced
Basis Methods and Reduction Error Models. JUQ 4(1):380--412.
\href{https://epubs.siam.org/doi/10.1137/140995817}{期刊页}，\href{https://eprints.whiterose.ac.uk/id/eprint/95248/11/WRRO_95248.pdf}{作者全文}。

{[}R12{]} Spantini A.; Cui T.; Willcox K.; Tenorio L.; Marzouk Y.
(2017). Goal-Oriented Optimal Approximations of Bayesian Linear Inverse
Problems. SISC 39(5):S167--S196.
\href{https://arxiv.org/abs/1607.01881}{作者稿}，10.1137/16M1082123。

{[}R13{]} Zahm O.; Cui T.; Law K.; Spantini A.; Marzouk Y. (2022).
Certified dimension reduction in nonlinear Bayesian inverse problems.
Math. Comp. 91:1789--1835.
\href{https://arxiv.org/abs/1807.03712}{作者稿}，10.1090/MCOM/3737。

{[}R14{]} Kaipio J. P.; Huttunen T.; Luostari T.; Lähivaara T.; Monk P.
B. (2019). A Bayesian approach to improving the Born approximation for
inverse scattering with high contrast materials. Inverse Problems.
\href{https://arxiv.org/abs/1901.00909}{作者稿}，10.1088/1361-6420/ab15f3。

{[}R15{]} Porat M.; Vinnikov V. (2021). Realizations of non-commutative
rational functions around a matrix centre, I. JLMS 104(3):1250--1299.
\href{https://londmathsoc.onlinelibrary.wiley.com/doi/10.1112/jlms.12459}{原文页}。

{[}R16{]} Kanagawa M.; Hennig P. (2019). Convergence Guarantees for
Adaptive Bayesian Quadrature Methods. NeurIPS.
\href{https://arxiv.org/abs/1905.10271}{作者稿}。

{[}R17{]} Wang Z. et al.~(2026-09-08). Finite-Modal Realization and
Operator-Norm Convergence of a Source-to-Observation Electromagnetic
Scattering Green Operator.
\href{https://arxiv.org/abs/2609.08229}{arXiv:2609.08229}。本轮只核到摘要，全文定理未核。

{[}R18{]} Geng N.; Jackson D. R.; Carin L. (1999). On the resonances of
a dielectric BOR buried in a dispersive layered medium. IEEE TAP
47(8):1305--1313.
\href{https://scholars.duke.edu/publication/749511}{作者机构记录}，10.1109/8.791948。

{[}R19{]} Wu T.; Jaramillo J. L.; Lalanne P. (2025). Reflections on the
Spatial Exponential Growth of Electromagnetic Quasinormal Modes. Laser
\& Photonics Reviews 19(9):2402133.
\href{https://doi.org/10.1002/lpor.202402133}{DOI}。用于限定复频环境敏感性，不将其结果当作本轮新定理。

{[}R20{]} Bui-Thanh T.; Ghattas O. (2013). Analysis of the Hessian for
inverse scattering problems. Part III: Inverse medium scattering of
electromagnetic waves in three dimensions. Inverse Problems and Imaging
7(4):1139--1155.
\href{https://www.aimsciences.org/article/doi/10.3934/ipi.2013.7.1139}{原文页}。

{[}R21{]} Li H.; Garg V. V.; Willcox K. (2018). Model adaptivity for
goal-oriented inference using adjoints. CMAME 331:1--22.
\href{https://dspace.mit.edu/entities/publication/af315104-eec6-4542-a3f4-bc37b9a3a16d}{作者机构原文}，10.1016/j.cma.2017.11.018。

{[}R22{]} Chen Q.; Arnaud É.; Baptista R.; Zahm O. (2025). Coupled
Input-Output Dimension Reduction: Application to Goal-oriented Bayesian
Experimental Design and Global Sensitivity Analysis. SISC
47(5):A2403--A2430.
\href{https://epubs.siam.org/doi/10.1137/24M1670421}{期刊页}。

附注：一般目标伴随、自适应模型精度和参数--观测联合降维也已有{[}R20--R22{]}等基础。上述推导应按具体Maxwell合同与新实验证据定位，不能仅因改用目标posterior语言就宣称原创。

\clearpage
\setlength{\parskip}{3pt}

\hypertarget{ux9644ux5f55bux4e0bux4e00ux8f6ecodexux6267ux884cux534fux8bae}{%
\section{附录B：下一轮Codex执行协议}\label{ux9644ux5f55bux4e0bux4e00ux8f6ecodexux6267ux884cux534fux8bae}}

完整独立文本随本报告提供。以下正文与独立Prompt一致；在执行时，以新目录保存结果，保持四个固定历史提交及其裁决不变。

Research W-GQ-R2 --- 下一轮 Codex 执行 Prompt

普通二维未知墙体下的目标分辨率、完整返回传播与目标充分概率计算

你负责在现有 migodam/b1-probabilistic-maxwell-imaging
仓库中执行一轮受控、可复现的研究。

同时提供的 \path{Research_W_GQ_R2_T_Report_ZH.pdf}
是理论研究，不是已运行实验。

目标是未知目标的真实误差、可辨识性、后验可靠性、完整Maxwell一致性与实际物理成本。

请在本合同和预算内自主完成所有可完成分支，不因一条算法失败终止其他已授权分支；不要仅写计划。

不得把任何Pro推导或预测记成Codex实测。不得修改旧提交、旧失败、旧账本或旧裁决。

\hypertarget{ux4e00ux5fc5ux987bux5148ux8bfbux7684ux56faux5b9aux6765ux6e90}{%
\subsection{一、必须先读的固定来源}\label{ux4e00ux5fc5ux987bux5148ux8bfbux7684ux56faux5b9aux6765ux6e90}}

\begin{enumerate}
\def\labelenumi{\arabic{enumi}.}
\tightlist
\item
  73ba6eec3746548bf21259f627b635678709df00:
\end{enumerate}

\path{research/research_w_bgga/RESEARCH_W_BGGA_REPORT_ZH.md}

\path{research/research_w_bgga/GATE_DECISION.json}

\path{research/research_w_bgga/COST_SUMMARY.json}

\path{research/research_w_bgga/CONTRACT.json}

及实际引用的G0/G1数组、查询接口、代价账本。

\begin{enumerate}
\def\labelenumi{\arabic{enumi}.}
\setcounter{enumi}{1}
\tightlist
\item
  6fd504785f4e4be74d919f2363e0b2d1ef111534:
\end{enumerate}

\path{research/research_w_bgga/RESEARCH_W_BGGA_REPORT_ZH.md}

注意这版是资格前检查，新增PDE为零。

\begin{enumerate}
\def\labelenumi{\arabic{enumi}.}
\setcounter{enumi}{2}
\tightlist
\item
  958a478739d9553269316ded8125fda9df225c0c:
\end{enumerate}

\path{research/research_w_p0r_to_p1/RESEARCH_REPORT_ZH.md}

\begin{enumerate}
\def\labelenumi{\arabic{enumi}.}
\setcounter{enumi}{3}
\tightlist
\item
  52351aee6ddd41b09851fe5e52d1a8861891ef13:
\end{enumerate}

\path{research/research_w_slam/RESEARCH_W_SLAM_REPORT_ZH.md}

research/research\_w\_slam/results/e2e3/finite\_bank/

\begin{enumerate}
\def\labelenumi{\arabic{enumi}.}
\setcounter{enumi}{4}
\tightlist
\item
  附件旧W-GQ与本轮R2理论报告。如有更新的Codex
  W-GQ账本，先定位它的精确提交与真实结果。
\end{enumerate}

不以当前main中未注明版本的文件替换冻结证据。

建新目录 research/research\_w\_gq\_r2/，保存：

\path{START_HERE.md}、\path{CONTRACT.json}、\path{PRE_REGISTERED_PLAN.md}、\path{THEORY_TO_TEST.md}、

\path{EVIDENCE_LEDGER.csv}、\path{QUERY_LEDGER.jsonl}、\path{PHYSICS_COST.json}、

\path{E1_RESULTS.csv}、\path{E2_RESULTS.csv}、\path{E3_RESULTS.csv}、

\path{THEOREM_STATUS.md}、\path{GO_NO_GO.json}、\path{REPORT_ZH.md}、\path{NEXT_PRO_HANDOFF.md}。

所有正文数字必须指向保存的实际数组或账本行。

\hypertarget{ux4e8cux4e0dux53efux6539ux53d8ux7684ux7269ux7406ux4e0eux63a8ux65adux5408ux540c}{%
\subsection{二、不可改变的物理与推断合同}\label{ux4e8cux4e0dux53efux6539ux53d8ux7684ux7269ux7406ux4e0eux63a8ux65adux5408ux540c}}

\begin{enumerate}
\def\labelenumi{\arabic{enumi}.}
\tightlist
\item
  二维非磁Ez scalar-TM，mu\_r=1，exp(-i omega t)，被动介质。
\end{enumerate}

主频率以f0和lambda0归一化；沿用已审核源、接收器、PML、单位线源、排序和白化。

\begin{enumerate}
\def\labelenumi{\arabic{enumi}.}
\setcounter{enumi}{1}
\tightlist
\item
  W、X、共享gamma在一组采集中静止。换目标配置是新独立世界，不是轨迹。
\end{enumerate}

solver迭代、换Green、查bank、算adjoint不是新观测。

\begin{enumerate}
\def\labelenumi{\arabic{enumi}.}
\setcounter{enumi}{2}
\tightlist
\item
  保持实虚每坐标sigma=0.01，不除通道数平方根。
\end{enumerate}

明确64实维基础数据与新增嵌套数据的维数。区分均值pair距离和含噪discrepancy。

\begin{enumerate}
\def\labelenumi{\arabic{enumi}.}
\setcounter{enumi}{3}
\tightlist
\item
  墙容器沿用BGGA公开ROI；目标ROI沿用公共目标域。
\end{enumerate}

真实mask、真实平均墙、真实目标位置不进入在线算法。

使用墙容器、已知均匀间隔、材料上下界或低维宏观参数时，要明确这是新受限合同。

\begin{enumerate}
\def\labelenumi{\arabic{enumi}.}
\setcounter{enumi}{4}
\tightlist
\item
  墙跨频共享epsilon(f)=epsilon\_inf+i
  ell/f；不允许每频独立造墙或任意gain。
\end{enumerate}

主仪器已校准；另设共享幅度{[}0.99,1.01{]}、相位{[}-0.01,0.01{]}压力测试。

\begin{enumerate}
\def\labelenumi{\arabic{enumi}.}
\setcounter{enumi}{5}
\tightlist
\item
  墙占据区域epsilon\_inf在{[}1.5,6{]}、ell在{[}0,0.25{]}；空气孔隙为epsilon=1。
\end{enumerate}

新微结构家族与旧minimum\_feature合同不同，要事前明确尺度，不能回写旧资格。

\begin{enumerate}
\def\labelenumi{\arabic{enumi}.}
\setcounter{enumi}{6}
\tightlist
\item
  新目标族：位置、尺寸、epsilon实部、损耗、圆/椭圆至少两分支。
\end{enumerate}

建议半径或等面积半径0.08/0.12/0.16 lambda0，epsilon实部2/4/7，ell
0/0.1/0.2；

椭圆纵横比1.5，固定两个事前方向。先小规模组合，不全笛卡尔膨胀。

若这些值与旧公共ROI有冲突，只在冻结前修正几何可行域，并在合同说明。

\begin{enumerate}
\def\labelenumi{\arabic{enumi}.}
\setcounter{enumi}{7}
\tightlist
\item
  散射场FEM必须保留材料依赖载荷：
\end{enumerate}

(L\_W - V\_X) u\_sc = b\_W + b\_X。

b\_X使用真实弱式incident积分，不能机械用节点incident或常右端替代。

\begin{enumerate}
\def\labelenumi{\arabic{enumi}.}
\setcounter{enumi}{8}
\tightlist
\item
  新传播表示必须保留目标自身与墙返回所有阶反馈：
\end{enumerate}

C\_W = C\_free + H\_W

T\_X\^{}0 = V\_X (I-C\_free V\_X)\^{}(-1)

F = d\_W + B\_W T\_X\^{}0 (I-H\_W T\_X\textsuperscript{0)}(-1) a\_W。

Born只作负对照，不作主物理替代。

\hypertarget{ux4e09ux8d44ux6e90ux4e0eux5ba1ux8ba1}{%
\subsection{三、资源与审计}\label{ux4e09ux8d44ux6e90ux4e0eux5ba1ux8ba1}}

本轮默认是小规模物理pilot，不启动大型PDE训练或全新巨大bank。

建议硬上限：

E1：2000 complex RHS，300次材料-频率FE分解，3 process CPU小时。

E2：2800 RHS，400 FE，3 CPU小时。

E3冷pilot：1200 RHS，200 FE，2 CPU小时。

总计6000 RHS、900 FE、8 CPU小时；阶段余额不自动转移。

峰值RSS不超过16GB与实际可用RAM的70\%两者较小值；实际机器不足则按比例缩小同时驻留。

E3旧bank逻辑开发不新增PDE，单独计CPU，不隐含免费teacher。

所有assembly、FE factorization、complex
RHS、adjoint、随机验证probe、失败、回退、

目标模板、dense闭环、离线分组及资格成本都入账。

wrapper CPU与内部kernel不重复相加；列出未计时缺口。

区分实测时间、工作量公式、历史已付成本、当前方法部署成本和评分reference成本。

达到上限即封账并报告未决，不松门槛、不隐藏失败、不自动扩为30-parent确认。

\hypertarget{ux56dbux9636ux6bb50ux6700ux5c0fux5fc5ux8981ux5ba1ux8ba1ux4e0eux96f6pdeux5f00ux53d1}{%
\subsection{四、阶段0：最小必要审计与零PDE开发}\label{ux56dbux9636ux6bb50ux6700ux5c0fux5fc5ux8981ux5ba1ux8ba1ux4e0eux96f6pdeux5f00ux53d1}}

\begin{enumerate}
\def\labelenumi{\arabic{enumi}.}
\tightlist
\item
  核对已有正负证据：
\end{enumerate}

G0精确吸收已通过；C2/C3不独立；C4目标不动；

M0位置MSE约0.35454 lambda0\^{}2；小regret不等于准确；

M4全查2950、M5保留118组；实际避免FE/RHS为0。

最终BGGA账本435.690883 CPU秒优先于正文早期421.3秒快照。

\begin{enumerate}
\def\labelenumi{\arabic{enumi}.}
\setcounter{enumi}{1}
\item
  阅读旧GQ报告，建立``旧已证明 / 本轮新推导 / 尚无实验''的表。
\item
  只为新接口做必要代数核对：
\end{enumerate}

T\_X\^{}0分离，contrast差分恒等式，四块余切，材料依赖载荷mixed
derivative。

不扩大无关unit tests，不重跑已合格的整套G0。

\begin{enumerate}
\def\labelenumi{\arabic{enumi}.}
\setcounter{enumi}{3}
\tightlist
\item
  冻结新world seeds、目标探针、质量门、预算、primary任务。
\end{enumerate}

所有后续改动登记原因；门槛不按结果移动。

\hypertarget{ux4e94e1ux7269ux7406ux5b58ux5728ux6027ux5728ux7ebfux8bc6ux522bux548cux6cdbux5316ux5206ux5f00}{%
\subsection{五、E1：物理存在性、在线识别和泛化分开}\label{ux4e94e1ux7269ux7406ux5b58ux5728ux6027ux5728ux7ebfux8bc6ux522bux548cux6cdbux5316ux5206ux5f00}}

\hypertarget{e1-pux4e3bux6b63ux4f8bux666eux901aux4e8cux7ef4ux5faeux7ed3ux6784}{%
\subsubsection{E1-P主正例：普通二维微结构。}\label{e1-pux4e3bux6b63ux4f8bux666eux901aux4e8cux7ef4ux5faeux7ed3ux6784}}

A. 选择至少两个非均匀宏观墙，含一个平滑或折线不规则边界。

B. 每个宏观墙构造两种局部复均值相同、纹理不同的有界普通材料。

可用二维颗粒与交错纹理；另有小空气孔隙臂，但保持合法体积分数与均值。

不把完全不可见区域当主机制。

C. 初始eta/lambda0建议1/8、1/12、1/16；以独立分辨率能解析为前提。

若更细网格成本不允许，标记尺度范围不足，不能用未解析微结构伪造等效。

D. 真实材料L2/RMS差应明显，建议RMS至少0.5；该值冻结后不追结果改。

E. 对一组事前冻结目标族计算完整响应，比较：

材料差、全任务最大白化响应差、闭环放大代理、目标与墙间距。

至少一个尺度达到pair距离\textless=1才记为噪声尺度等效正例；

同时保留未达1但呈尺度趋势的机制证据，两者不能混为成功。

F. 至少一个强散射目标要求Born与full明显不同。

independent mesh/PML/积分检查用于最差探针，不能借另一个probe的误差。

经验网格差不叫连续统一证书。

\hypertarget{e1-sux7ed3ux6784ux6b63ux5bf9ux7167ux955cux50cfux5899}{%
\subsubsection{E1-S结构正对照：镜像墙。}\label{e1-sux7ed3ux6784ux6b63ux5bf9ux7167ux955cux50cfux5899}}

A. 任意二维镜像非均匀墙，允许对称空腔，目标居中。

B. 成对源/收与一致白化；检验P\_- A近零、P\_+ B\_x近零及B\_x非零。

C. 加入非对称墙扰动或左右仪器偏差，检查保护是否丧失。

D. 只宣称居中横向任务的结论；离中心、深度、材料估计另评。

\hypertarget{e1-nux8d1fux4f8bux5168ux6ce2ux4ea7ux54c1ux6838ux6ee1ux5c04ux4e0eux6709ux9650ux8865ux507f}{%
\subsubsection{E1-N负例：全波产品核满射与有限补偿。}\label{e1-nux8d1fux4f8bux5168ux6ce2ux4ea7ux54c1ux6838ux6ee1ux5c04ux4e0eux6709ux9650ux8865ux507f}}

A. 在严格被动墙内选允许的epsilon/loss bump，组装实白化墙Jacobian。

B. 报告秩、最小奇异值、导数数值误差；不得把小非零奇异值随意删掉。

C.
若确有满射裕度，指定目标方向，用有界continuation+Full-J正确器追踪补偿。

D. 验证实际full响应、合法性、独立网格、有限位移尺度和新增观测破坏。

E. 未找到分支标为优化/预算未达；不判不存在多解。

\hypertarget{e1-oux5728ux7ebfux672aux77e5ux76eeux6807-1}{%
\subsubsection{E1-O在线未知目标。}\label{e1-oux5728ux7ebfux672aux77e5ux76eeux6807-1}}

A. 另冻结独立位置、半径、材料、损耗和形状parent，建议初始12
parent、每parent2噪声。

B. 数据生成进程保存truth；在线只读y、prior和公开几何容器。

C. 比较：

真实墙oracle（诊断上界）；

zero-data prior；

充分的普通联合Full-J/TR或L-BFGS/多起点FWI；

普通先估墙后固定；

未知宏观参数的有效传播推断；

保留多背景的目标marginal方法。

D. 不用真实均值墙作在线初值；不从truth筛选正确branch。

支撑搜索与连续优化分开归因。

给普通基线充分同预算，记录stationarity/约束残差；不能用旧6次短拟合当成熟FWI。

E. 主定位门在新数据打开前冻结为RMSE\textless=0.1
lambda0，同时报中位、90分位、最坏误差。

半径、材料、损耗和形状分别评分；不把位置成功写成所有属性准确。

F. 输出先封存再评分。若仅oracle成功，标为存在性成功、在线识别未建立。

\hypertarget{e1-gux6cdbux5316ux4e0eux7834ux574f}{%
\subsubsection{E1-G泛化与破坏。}\label{e1-gux6cdbux5316ux4e0eux7834ux574f}}

冻结未用于构造的更近墙、更大目标、更高频、另一形状、低损耗强反馈与宏观空腔。

同一墙对不能针对每个测试目标再拟合。

分别诊断尺度失效、反馈共振、宏观模型错误、观测不足与优化未达。

\hypertarget{ux516de2ux540cux4e00ux4e16ux754cux7684ux771fux6b63ux5d4cux5957ux6d4bux91cf}{%
\subsection{六、E2：同一世界的真正嵌套测量}\label{ux516de2ux540cux4e00ux4e16ux754cux7684ux771fux6b63ux5d4cux5957ux6d4bux91cf}}

A0：原两发16收、f0。

A1：完整保留A0数据，加入原外侧两个源、同16收、f0。

A2：保留A1，加入0.85 f0与1.15 f0的四源16收。

新noise独立生成，旧noise保留；固定总能量导致旧SNR改变的设计另列，不能套嵌套定理。

同时测：

\begin{enumerate}
\def\labelenumi{\arabic{enumi}.}
\item
  K\_W\textbar X,gamma 与
  K\_X\textbar W,gamma，位置/尺寸/材料任务分别尺度化。
\item
  exact algebraic rank、数值阈值rank、物理幅度受限信息分列。
\end{enumerate}

A\_full可能满行而K\_exact=0；截小mode后信息必须标constraint/noise
contract。

\begin{enumerate}
\def\labelenumi{\arabic{enumi}.}
\setcounter{enumi}{2}
\tightlist
\item
  多个目标方向的有限profile：
\end{enumerate}

给定目标位移，优化合法W及其他nuisance；

保存全波端点、实际材料变动、误差、convergence、budget。

只有找到的pair能给omega下界；没找到不是omega上界。

\begin{enumerate}
\def\labelenumi{\arabic{enumi}.}
\setcounter{enumi}{3}
\tightlist
\item
  二阶mixed compensation：
\end{enumerate}

不只用一个A伪逆解，要允许ker(A)及合法边界方向；

支撑/拓扑间用finite secant。

\begin{enumerate}
\def\labelenumi{\arabic{enumi}.}
\setcounter{enumi}{4}
\tightlist
\item
  若尝试有限盒上界，验证J0伪逆下的Jacobian变化界\textless1；
\end{enumerate}

抽样最大值只是经验界，不能当整个盒的认证sup。

\begin{enumerate}
\def\labelenumi{\arabic{enumi}.}
\setcounter{enumi}{5}
\item
  目标真实误差、posterior风险、CDF/event、形状概率、coverage与predictive
  score。
\item
  同时给wall材料方差/误差，寻找目标明显改善而wall细节仍宽的世界。
\item
  新测量单次真值误差可以上升；按parent配对汇总，不删反例。
\item
  不以新的source/frequency测量成本冒充内部solver迭代成本。
\end{enumerate}

\hypertarget{ux4e03e3ux5148ux65e7bankux540eux6709ux9650ux51b7ux542fux52a8}{%
\subsection{七、E3：先旧bank，后有限冷启动}\label{ux4e03e3ux5148ux65e7bankux540eux6709ux9650ux51b7ux542fux52a8}}

\hypertarget{e3-lux53d7ux63a7ux67e5ux8be2ux63a5ux53e3}{%
\subsubsection{E3-L受控查询接口。}\label{e3-lux53d7ux63a7ux67e5ux8be2ux63a5ux53e3}}

\begin{enumerate}
\def\labelenumi{\arabic{enumi}.}
\tightlist
\item
  旧118墙×25位置、2950原子仅为明确离散counting prior。
\end{enumerate}

target intrinsic/gamma已知；不能声称验证未知材料、孔洞或连续posterior。

\begin{enumerate}
\def\labelenumi{\arabic{enumi}.}
\setcounter{enumi}{1}
\tightlist
\item
  在线进程仅见prior元数据、主观测y和query(atom\_id)。
\end{enumerate}

外侧TX2/3、全部表、完整posterior、truth禁止预读。

每个响应、导数、离线teacher和代表模型构建都记录。

\begin{enumerate}
\def\labelenumi{\arabic{enumi}.}
\setcounter{enumi}{2}
\tightlist
\item
  比较完整response距离、目标increment距离、
\end{enumerate}

包含背景残差投影的contrast距离、条件目标posterior/decision距离。

\begin{enumerate}
\def\labelenumi{\arabic{enumi}.}
\setcounter{enumi}{3}
\tightlist
\item
  不免费删除每背景anchor likelihood：
\end{enumerate}

D\_W(X)=\textless r\_W0,Delta\_WX\textgreater-0.5\textbar\textbar Delta\_WX\textbar\textbar\^{}2。

相同Delta不够，还要控制背景残差投影。

\begin{enumerate}
\def\labelenumi{\arabic{enumi}.}
\setcounter{enumi}{4}
\tightlist
\item
  新机制是conditional consensus：
\end{enumerate}

p\_X=sum beta\_g p\_g。

若所有p\_g距代表\textless=epsilon，可以不准确计算beta；

若仅共同动作regret小，只输出decision certificate；

若目标分歧大，必须计算组间质量。

相关prior需包括conditional target prior；非法anchor用共同支撑分块。

\begin{enumerate}
\def\labelenumi{\arabic{enumi}.}
\setcounter{enumi}{5}
\tightlist
\item
  实现基线：
\end{enumerate}

全枚举；

普通自适应求积/默认未知质量包络；

普通response ROM/插值；

普通分离返回表示+统一输出精度；

同表示+contrast余切/二阶界+conditional consensus。

\begin{enumerate}
\def\labelenumi{\arabic{enumi}.}
\setcounter{enumi}{6}
\tightlist
\item
  任何从完整表离线得到的分组/谱/方差仅作offline诊断，不记在线节省。
\end{enumerate}

如果现有bank缺少物理证书，不得用完整表拟合证书后当零成本oracle。

\hypertarget{e3-pux6709ux9650ux51b7pdeux4ec5ux5728ux4e0aux8ff0ux673aux5236ux6709ux903bux8f91ux4e0eux53efux83b7ux5f97ux8bc1ux4e66ux8bc1ux636eux540e}{%
\subsubsection{E3-P有限冷PDE（仅在上述机制有逻辑与可获得证书证据后）。}\label{e3-pux6709ux9650ux51b7pdeux4ec5ux5728ux4e0aux8ff0ux673aux5236ux6709ux903bux8f91ux4e0eux53efux83b7ux5f97ux8bc1ux4e66ux8bc1ux636eux540e}}

\begin{enumerate}
\def\labelenumi{\arabic{enumi}.}
\item
  新小prior建议8墙×25目标原子，至少两种内禀模板；先冻结再生成观测。
\item
  保留C\_free精确自散射，只压缩H\_wall；r从4、8、16、32逐步提升。
\end{enumerate}

不把r等同旧22列goal-ROM；旧失败不直接决定新表示。

\begin{enumerate}
\def\labelenumi{\arabic{enumi}.}
\setcounter{enumi}{2}
\tightlist
\item
  固定圆盘位置变化可共享每频Mie/T模板及平移；普通基线同享。
\end{enumerate}

FEM/PML与连续自由空间模板混接必须独立检查。

\begin{enumerate}
\def\labelenumi{\arabic{enumi}.}
\setcounter{enumi}{3}
\item
  用普通全域Full、普通Schur/两物体T-matrix、Krylov、goal-ROM和概率策略作同质量对照。
\item
  若用随机action证书：
\end{enumerate}

E固定于此前历史，fresh g\textasciitilde N(0,I)覆盖完整声明离散输入空间；

t=Phi\textsuperscript{-1((1+alpha}(1/q))/2)；

\textbar\textbar E\textbar\textbar\textless=max\textbar\textbar Eg\textbar\textbar/t以概率\textgreater=1-alpha成立。

计入求解误差，使用正确迹范数，每wall/frequency/轮次做alpha spending。

只认证离散E，不能覆盖未建接口尾、连续PDE和PML。

alpha\_total建议1e-4，事前冻结；这是算法验证随机性而非Bayes可信度。

\begin{enumerate}
\def\labelenumi{\arabic{enumi}.}
\setcounter{enumi}{5}
\item
  闭环稳定界与contrast二阶余项必须使用；当前一阶score近零不够。
\item
  每种部署从相同空缓存、同权限启动，实际FE/RHS/CPU/RSS全部计入。
\item
  查询数、RHS数、FE数和总秒数分别判门；不得用较少表访问代替冷部署加速。
\item
  若证书算得比Full贵，记录CERTIFICATE\_COST\_FAILURE；不降低质量标准。
\end{enumerate}

\hypertarget{ux516bux51bbux7ed3ux8d28ux91cfux95e8ux4e0eux6982ux7387ux8303ux56f4}{%
\subsection{八、冻结质量门与概率范围}\label{ux516bux51bbux7ed3ux8d28ux91cfux95e8ux4e0eux6982ux7387ux8303ux56f4}}

\begin{enumerate}
\def\labelenumi{\arabic{enumi}.}
\item
  物理：主检验保留完整目标反馈和材料右端；每个关键结论有独立数值资格。
\item
  定位：新E1/E2主任务RMSE\textless=0.1
  lambda0；给与prior及oracle的差距。
\end{enumerate}

如果只某个受限对称点通过，不能扩成一般二维定位。

\begin{enumerate}
\def\labelenumi{\arabic{enumi}.}
\setcounter{enumi}{2}
\item
  目标posterior：有限reference下TV、事件/CDF误差建议\textless=0.02。
\item
  决策：平方位置regret另列，绝不替代绝对MSE。
\item
  coverage：给parent-level区间，12
  parent只是pilot；不以24/24覆盖宣布精确校准。
\item
  posterior数值界、FEM误差、模型误差分开。
\end{enumerate}

若osc log(mtilde/m)\textless=kappa，可用tanh(kappa/4)；

forward白化0.1不自动认证TV0.02。

\begin{enumerate}
\def\labelenumi{\arabic{enumi}.}
\setcounter{enumi}{6}
\tightlist
\item
  保留浮点上下界正越界，不用新增容差抹除。
\end{enumerate}

对固定浮点log-potential的高精度归一化证书不升级为物理posterior证书。

\begin{enumerate}
\def\labelenumi{\arabic{enumi}.}
\setcounter{enumi}{7}
\item
  所有数据分支、形状分支及未计算质量有记录；ESS不能认证遗漏质量。
\item
  新模型包含的microstructure平均机制是受限物理合同，不是任意墙正确的证明。
\item
  无新增NN、3D、大模型训练或无关A17模块。
\end{enumerate}

\hypertarget{ux4e5dux5fc5ux987bux8f93ux51faux7684ux79d1ux5b66ux56feux4e0eux8868}{%
\subsection{九、必须输出的科学图与表}\label{ux4e5dux5fc5ux987bux8f93ux51faux7684ux79d1ux5b66ux56feux4e0eux8868}}

图1：同均值不同二维墙材料；对应整个冻结目标族响应差与eta。

图2：真实目标误差与wall误差/不确定性二维散点，标prior和oracle。

图3：嵌套测量下目标绝对信息、有限补偿、目标风险与wall剩余不确定性。

图4：Born、完整目标自散射、完整墙返回反馈的受控差异。

图5：完整response、increment、contrast、conditional-target分歧的比较。

图6：相同质量下query/RHS/FE/CPU曲线，setup/证书/fallback全列。

图7：一个按理论发生的失效反例（近墙、非对称、共振或非法gauge）。

这些图只能使用保存的实际数据；不手绘预期``成像成功''图。

统计先在同parent内汇总重复noise，再进行paired分析，报告全部parent与失败分母。

\hypertarget{ux5341ux6700ux7ec8ux4ea4ux4ed8ux4e0eux5206ux5408ux540cux88c1ux51b3}{%
\subsection{十、最终交付与分合同裁决}\label{ux5341ux6700ux7ec8ux4ea4ux4ed8ux4e0eux5206ux5408ux540cux88c1ux51b3}}

每个命题标：

PROVED\_UNDER\_STATED\_ASSUMPTIONS / NUMERICALLY\_SUPPORTED /
NOT\_ESTABLISHED /

OPTIMIZATION\_NOT\_REACHED / MODEL\_MISMATCH / IDENTIFIABILITY\_FAILURE
/

COMPRESSION\_FAILURE / CERTIFICATE\_COST\_FAILURE /
NET\_SPEEDUP\_NOT\_ESTABLISHED。

至少回答：

A. 普通二维wrong wall--correct target是否有新完整全波正例？

B. 只是存在性，还是measurement-only在线成功？

C. 哪些目标/频率/孔径扩展破坏原等价？

D. 目标改善是否同时要求wall材料改善？

E. 是否必须精算所有背景质量？

F. 哪个具体Green action或FE/RHS被安全跳过？

G. 收益来自标准T-matrix/Schur结构还是新增目标概率规则？

H. 目标绝对误差是否真正下降，而非仅复制不准确reference？

I. 连续Maxwell、全prior和一般unknown-support哪些仍未认证？

J. 哪一条最值得形成IEEE TAP贡献，哪一条应停止？

最终提供完整中文报告、机器判决、可复现命令、全部成本与失败、

关键实际数组、数学命题的实现合同以及下一轮Pro handoff。

不把一次局部优化失败当成无解；不因一个算法不加速终结目标可辨识性主线。

\end{document}
